\documentclass[aps,preprint]{revtex4}%
\usepackage{amsfonts}
\usepackage{amsmath}
\usepackage{amssymb}
\usepackage{graphicx}%
\setcounter{MaxMatrixCols}{30}
%TCIDATA{OutputFilter=latex2.dll}
%TCIDATA{Version=4.10.0.2363}
%TCIDATA{CSTFile=revtex4.cst}
%TCIDATA{Created=Sunday, April 20, 2003 13:13:14}
%TCIDATA{LastRevised=Friday, May 09, 2003 10:49:01}
%TCIDATA{<META NAME="GraphicsSave" CONTENT="32">}
%TCIDATA{<META NAME="DocumentShell" CONTENT="Articles\SW\REVTeX 4">}
\newtheorem{theorem}{Theorem}
\newtheorem{acknowledgement}[theorem]{Acknowledgement}
\newtheorem{algorithm}[theorem]{Algorithm}
\newtheorem{axiom}[theorem]{Axiom}
\newtheorem{claim}[theorem]{Claim}
\newtheorem{conclusion}[theorem]{Conclusion}
\newtheorem{condition}[theorem]{Condition}
\newtheorem{conjecture}[theorem]{Conjecture}
\newtheorem{corollary}[theorem]{Corollary}
\newtheorem{criterion}[theorem]{Criterion}
\newtheorem{definition}[theorem]{Definition}
\newtheorem{example}[theorem]{Example}
\newtheorem{exercise}[theorem]{Exercise}
\newtheorem{lemma}[theorem]{Lemma}
\newtheorem{notation}[theorem]{Notation}
\newtheorem{problem}[theorem]{Problem}
\newtheorem{proposition}[theorem]{Proposition}
\newtheorem{remark}[theorem]{Remark}
\newtheorem{solution}[theorem]{Solution}
\newtheorem{summary}[theorem]{Summary}
\newenvironment{proof}[1][Proof]{\noindent\textbf{#1.} }{\ \rule{0.5em}{0.5em}}
\begin{document}
\preprint{HEP/123-qed}
\title[Short title for running header]{Long title that appears on the title page}
\author{First Author}
\affiliation{My Institution}
\author{Second Author}
\affiliation{My Institution}
\author{Third Author}
\affiliation{Other Institution}
\keywords{one two three}
\pacs{PACS number}

\begin{abstract}
Shell document for REV\TeX{} 4.

\end{abstract}
\volumeyear{year}
\volumenumber{number}
\issuenumber{number}
\eid{identifier}
\date[Date text]{date}
\received[Received text]{date}

\revised[Revised text]{date}

\accepted[Accepted text]{date}

\published[Published text]{date}

\startpage{101}
\endpage{102}
\maketitle
\tableofcontents


\subsubsection{Generalized Bases for $\mathcal{B}\left(  \mathcal{H}%
^{N}\right)  $ and $\mathcal{B}_{1}\left(  \mathcal{H}^{N}\right)  $}

The duality between $\mathcal{B}_{1}\left(  \mathcal{H}^{N}\right)  $ and
$\mathcal{B}\left(  \mathcal{H}^{N}\right)  $ can be expressed in terms of the
scalar product
\begin{equation}
\left(  X|Y\right)  =Tr\left\{  X^{\dagger}Y\right\}  ;\quad X\in
\mathcal{B}\left(  \mathcal{H}^{N}\right)  ,Y\in\mathcal{B}_{1}\left(
\mathcal{H}^{N}\right)
\end{equation}
Following the Dirac convention we can denote operators, $Y\in\mathcal{B}%
_{1}\left(  \mathcal{H}^{N}\right)  $, as $\left\vert Y\right)  $ and
operators, $X\in\mathcal{B}\left(  \mathcal{H}^{N}\right)  $, as $\left(
X\right\vert $ thus producing the super operator%
\begin{equation}
\left\vert Y\right)  \left(  X\right\vert :\mathcal{B}_{1}\left(
\mathcal{H}^{N}\right)  \rightarrow\mathcal{B}_{1}\left(  \mathcal{H}%
^{N}\right)
\end{equation}
whose action is given by
\begin{equation}
\left\vert Y\right)  \left(  X\right\vert \,Z=Tr\left\{  X^{\dagger}Z\right\}
\left\vert Y\right)  \equiv Tr\left\{  X^{\dagger}Z\right\}  Y
\end{equation}


We can introduce a generalized basis, $\left\{  \left\vert \mathsf{z}%
^{N}\right\rangle \left\langle \mathsf{z}^{\prime N}\right\vert \right\}  $,
\emph{for both} of the dual spaces $\mathcal{B}\left(  \mathcal{H}^{N}\right)
$ and $\mathcal{B}_{1}\left(  \mathcal{H}^{N}\right)  $, composed of ket-bra
operators that correspond to the generalized eigenvectors of the $N$-particle
space-spin coordinate operator. None of the basis elements $\left\vert
\mathsf{z}^{N}\right\rangle \left\langle \mathsf{z}^{\prime N}\right\vert $
belong to $\mathcal{B}\left(  \mathcal{H}^{N}\right)  $ as the action of any
individual $\left\vert \mathsf{z}^{N}\right\rangle \left\langle \mathsf{z}%
^{\prime N}\right\vert $ on any vector belonging to $\mathcal{H}^{N}$ produces
a vector $\notin\mathcal{H}^{N}$. The notation $\left\vert \left\vert
\mathsf{z}^{N}\right\rangle \left\langle \mathsf{z}^{\prime N}\right\vert
\right)  $ signifies that $\left\vert \mathsf{z}^{N}\right\rangle \left\langle
\mathsf{z}^{\prime N}\right\vert $ is to be considered within the
$\mathcal{B}_{1}\left(  \mathcal{H}^{N}\right)  $-topology, while the notation
$\left(  \left\vert \mathsf{z}^{N}\right\rangle \left\langle \mathsf{z}%
^{\prime N}\right\vert \right\vert $ signifies that $\left\vert \mathsf{z}%
^{N}\right\rangle \left\langle \mathsf{z}^{\prime N}\right\vert $ is to be
considered within the $\mathcal{B}\left(  \mathcal{H}^{N}\right)  $-topology.

This basis allows one to define a super operator resolution of the identity on
$\mathcal{B}_{1}\left(  \mathcal{H}^{N}\right)  $ by
\begin{equation}
\hat{I}_{\mathcal{B}_{1}\left(  \mathcal{H}^{N}\right)  }=%
%TCIMACRO{\dint }%
%BeginExpansion
{\displaystyle\int}
%EndExpansion
\left\vert \left\vert \mathsf{z}^{N}\right\rangle \left\langle \mathsf{z}%
^{\prime N}\right\vert \right)  \left(  \left\vert \mathsf{z}^{N}\right\rangle
\left\langle \mathsf{z}^{\prime N}\right\vert \right\vert d\mathsf{z}%
^{N}d\mathsf{z}^{\prime N},
\end{equation}
which can be expressed symbolically as
\begin{equation}
\hat{I}_{\mathcal{B}_{1}\left(  \mathcal{H}^{N}\right)  }=\left(  \left(
\underline{\left\vert \mathsf{z}^{N}\right\rangle \left\langle \mathsf{z}%
^{\prime N}\right\vert }\right\vert \right)  \left(  \left\vert \underline
{\left\vert \mathsf{z}^{N}\right\rangle \left\langle \mathsf{z}^{\prime
N}\right\vert }\right)  \right)  ^{t},
\end{equation}
producing the generalized expansion of an element $X\in\mathcal{B}_{1}\left(
\mathcal{H}^{N}\right)  $ as%
\begin{align}
X  &  =%
%TCIMACRO{\dint }%
%BeginExpansion
{\displaystyle\int}
%EndExpansion
\left\vert \left\vert \mathsf{z}^{N}\right\rangle \left\langle \mathsf{z}%
^{\prime N}\right\vert \right)  \left(  \left\vert \mathsf{z}^{N}\right\rangle
\left\langle \mathsf{z}^{\prime N}\right\vert |X\right)  d\mathsf{z}%
^{N}d\mathsf{z}^{\prime N}\nonumber\\
&  =%
%TCIMACRO{\dint }%
%BeginExpansion
{\displaystyle\int}
%EndExpansion
\left\vert \left\vert \mathsf{z}^{N}\right\rangle \left\langle \mathsf{z}%
^{\prime N}\right\vert \right)  Tr\left\{  \left\vert \mathsf{z}%
^{N}\right\rangle \left\langle \mathsf{z}^{\prime N}\right\vert ^{\dagger
}X\right\}  d\mathsf{z}^{N}d\mathsf{z}^{\prime N}\nonumber\\
&  =%
%TCIMACRO{\dint }%
%BeginExpansion
{\displaystyle\int}
%EndExpansion
\left\vert \left\vert \mathsf{z}^{N}\right\rangle \left\langle \mathsf{z}%
^{\prime N}\right\vert \right)  \left\langle \mathsf{z}^{N}|X\mathsf{z}%
^{\prime N}\right\rangle d\mathsf{z}^{N}d\mathsf{z}^{\prime N}\nonumber\\
&  =%
%TCIMACRO{\dint }%
%BeginExpansion
{\displaystyle\int}
%EndExpansion
X\left(  \mathsf{z}^{N},\mathsf{z}^{\prime N}\right)  \left\vert
\mathsf{z}^{N}\right\rangle \left\langle \mathsf{z}^{\prime N}\right\vert
d\mathsf{z}^{N}d\mathsf{z}^{\prime N}%
\end{align}
where $X\left(  \mathsf{z}^{N},\mathsf{z}^{\prime N}\right)  $ is the integral
kernel corresponding to $X$.

The basis $\left\{  \left\vert \mathsf{z}^{N}\right\rangle \left\langle
\mathsf{z}^{\prime N}\right\vert \right\}  $ can be graded into sub bases
$\left\{  \mathfrak{b}_{m}\left(  \mathsf{z}\right)  \right\}  $ with $m$ bra
and ket coordinates equal which are defined as
\begin{equation}
\mathfrak{b}_{m}\left(  \mathsf{z}\right)  =\left\{  \left\vert \mathsf{z}%
^{N}\right\rangle \left\langle \mathsf{z}^{\prime N}\right\vert |\mathsf{z}%
^{N}\overset{m}{=}\mathsf{z}^{\prime N}\right\}
\end{equation}
where $\mathsf{z}^{N}\overset{m}{=}\mathsf{z}^{\prime N}$ denotes that the
coordinates $\mathsf{z}^{N}$ and $\mathsf{z}^{\prime N}$ have exactly $m$
values in common, which could be equivalently written as $\mathsf{z}%
^{N}\overset{N-m}{\neq}\mathsf{z}^{\prime N}$ i.e. $\mathsf{z}^{N}$ and
$\mathsf{z}^{\prime N}$ have exactly $N-m$ values different. These bases
produce a graded decomposition of $\mathcal{B}_{1}\left(  \mathcal{H}%
^{N}\right)  $ as the direct sum
\begin{equation}
\mathcal{B}_{1}\left(  \mathcal{H}^{N}\right)  =%
%TCIMACRO{\dbigoplus \limits_{0\leq m\leq N}}%
%BeginExpansion
{\displaystyle\bigoplus\limits_{0\leq m\leq N}}
%EndExpansion
\mathfrak{B}_{m}\left(  \mathsf{z}\right)
\end{equation}
where $\mathfrak{B}_{m}\left(  \mathsf{z}\right)  $ is defined as the closed
linear span of the basis $\left\{  \mathfrak{b}_{m}\left(  \mathsf{z}\right)
\right\}  $ i.e.
\begin{equation}
\mathfrak{B}_{m}\left(  \mathsf{z}\right)  =\overline{\operatorname{linspan}%
}\left\{  \left.  \underset{}{\left\vert \mathsf{z}^{N}\right\rangle
\left\langle \mathsf{z}^{\prime N}\right\vert }\right\vert \left\vert
\mathsf{z}^{N}\right\rangle \left\langle \mathsf{z}^{\prime N}\right\vert
\in\mathfrak{b}_{m}\left(  \mathsf{z}\right)  \right\}
\end{equation}
and the closure is taken wrt the trace norm topology. In order to analyze the
kernels of the contraction maps $\left\{  L_{N}^{p}|0\leq p\leq N\right\}  $
it is useful to consider the sequence of subspaces $\left\{  \mathfrak{V}%
_{p}\left(  \mathsf{z}\right)  \right\}  $, that are generated by $\left\{
\left\vert \mathsf{z}^{N}\right\rangle \left\langle \mathsf{z}^{\prime
N}\right\vert \right\}  $ that have at least $p$ differences between
$\mathsf{z}^{N}$ and $\mathsf{z}^{\prime N}$ (denoted by $\mathsf{z}%
^{N}\overset{p}{\nsim}$ $\mathsf{z}^{\prime N}$) or equivalently generated by
bases elements that have at most $N-p$ values in common between $\mathsf{z}%
^{N}$ and $\mathsf{z}^{\prime N}$ (denoted by $\mathsf{z}^{N}\overset
{N-p}{\sim}$ $\mathsf{z}^{\prime N})$, which are defined as
\begin{equation}
\mathfrak{V}_{p}\left(  \mathsf{z}\right)  =%
%TCIMACRO{\dbigoplus \limits_{0\leq m\leq N-p}}%
%BeginExpansion
{\displaystyle\bigoplus\limits_{0\leq m\leq N-p}}
%EndExpansion
\mathfrak{B}_{m}\left(  \mathsf{z}\right)
\end{equation}
so that
\begin{equation}
\mathfrak{V}_{p}\left(  \mathsf{z}\right)  =\overline{\operatorname{linspan}%
}\left\{  \left.  \left\vert \mathsf{z}^{N}\right\rangle \left\langle
\mathsf{z}^{\prime N}\right\vert \overset{}{}\right\vert \mathsf{z}%
^{N}\overset{p}{\nsim}\mathsf{z}^{\prime N}\right\}  =\overline
{\operatorname{linspan}}\left\{  \mathfrak{c}_{p}\left(  \mathsf{z}\right)
\right\}
\end{equation}
where $\mathfrak{c}_{p}\left(  \mathsf{z}\right)  \equiv\left\{  \left.
\underset{}{\left\vert \mathsf{z}^{N}\right\rangle \left\langle \mathsf{z}%
^{\prime N}\right\vert }\right\vert \mathsf{z}^{N}\overset{p}{\nsim}%
\mathsf{z}^{\prime N}\right\}  $ $\equiv\left\{  \left.  \underset
{}{\left\vert \mathsf{z}^{N}\right\rangle \left\langle \mathsf{z}^{\prime
N}\right\vert }\right\vert \mathsf{z}^{N}\overset{N-p}{\sim}\mathsf{z}^{\prime
N}\right\}  $ and hence $\mathfrak{V}_{0}\left(  \mathsf{z}\right)
=\mathcal{B}_{1}\left(  \mathcal{H}^{N}\right)  $. One should also observe
that
\begin{equation}
\mathfrak{B}_{N-p}\left(  \mathsf{z}\right)  =\mathfrak{V}_{p-1}\left(
\mathsf{z}\right)  \ominus\mathfrak{V}_{p}\left(  \mathsf{z}\right)
\end{equation}
These subspaces partially characterize the kernel of the maps $\left\{
L_{N}^{p}\right\}  $ as they have the property
\begin{equation}
\mathfrak{V}_{p-1}\left(  \mathsf{z}\right)  \supset\ker L_{N}^{p}%
\supset\mathfrak{V}_{p}\left(  \mathsf{z}\right)  .
\end{equation}
and one has in general%
\begin{equation}
\mathcal{B}_{1}\left(  \mathcal{H}^{N}\right)  =\mathfrak{V}_{0}\left(
\mathsf{z}\right)  \supset\ker L_{N}^{0}\supset\mathfrak{V}_{1}\left(
\mathsf{z}\right)  \supset\ker L_{N}^{1}\supset\mathfrak{V}_{2}\left(
\mathsf{z}\right)  \supset\cdots\supset\mathfrak{V}_{N}\left(  \mathsf{z}%
\right)  \supset\ker L_{N}^{N}=\phi
\end{equation}


\section{Proper Bases for $\mathcal{B}\left(  \mathcal{H}^{N}\right)  $ and
$\mathcal{B}_{1}\left(  \mathcal{H}^{N}\right)  $}

One can repeat all the above for proper bases of $\mathcal{H}^{N}$ generated
by complete orthonormal bases $\left\{  \varphi_{j}\right\}  $ of
$\mathcal{H}^{N}$ to produce $\mathfrak{B}_{p}\left(  \mathbf{\varphi}\right)
,$ $\mathfrak{b}_{p}\left(  \mathbf{\varphi}\right)  ,\mathfrak{c}_{p}\left(
\mathbf{\varphi}\right)  $ and $\mathfrak{V}_{p}\left(  \mathbf{\varphi
}\right)  $ and their associated decompositions of $\mathcal{B}_{1}\left(
\mathcal{H}^{N}\right)  .$ Hence in particular
\begin{equation}
\mathfrak{V}_{p-1}\left(  \mathbf{\varphi}\right)  \supset\ker L_{N}%
^{p}\supset\mathfrak{V}_{p}\left(  \mathbf{\varphi}\right)
\end{equation}%
\begin{equation}
\mathfrak{V}_{p}\left(  \mathbf{\varphi}\right)  =%
%TCIMACRO{\dbigoplus \limits_{0\leq m\leq N-p}}%
%BeginExpansion
{\displaystyle\bigoplus\limits_{0\leq m\leq N-p}}
%EndExpansion
\mathfrak{B}_{m}\left(  \mathbf{\varphi}\right)
\end{equation}%
\begin{equation}
\mathfrak{B}_{m}\left(  \mathbf{\varphi}\right)  =\overline
{\operatorname{linspan}}\left\{  \left.  \underset{}{\left\vert
\mathbf{\varphi}^{N}\right\rangle \left\langle \mathbf{\varphi}^{\not \prime
N}\right\vert }\right\vert \left\vert \mathbf{\varphi}^{N}\right\rangle
\left\langle \mathbf{\varphi}^{\not \prime N}\right\vert \in\mathfrak{b}%
_{m}\left(  \mathbf{\varphi}\right)  \right\}
\end{equation}%
\begin{equation}
\mathfrak{V}_{p}\left(  \mathbf{\varphi}\right)  =\overline
{\operatorname{linspan}}\left\{  \left.  \left\vert \mathbf{\varphi}%
^{N}\right\rangle \left\langle \mathbf{\varphi}^{\not \prime N}\right\vert
\overset{}{}\right\vert \mathbf{\varphi}^{N}\overset{p}{\nsim}\mathbf{\varphi
}^{\not \prime N}\right\}  =\overline{\operatorname{linspan}}\left\{
\mathfrak{c}_{p}\left(  \mathbf{\varphi}\right)  \right\}
\end{equation}%
\begin{equation}
\mathfrak{B}_{N-p}\left(  \mathbf{\varphi}\right)  =\mathfrak{V}_{p-1}\left(
\mathbf{\varphi}\right)  \ominus\mathfrak{V}_{p}\left(  \mathbf{\varphi
}\right)
\end{equation}


\section{Representations of $U\left(  \mathcal{H}^{1}\right)  $}

One can define subspaces of $\mathcal{B}_{1}\left(  \mathcal{H}^{N}\right)  $
that produce irreducible representations of $SU\left(  \mathcal{H}^{1}\right)
$ that are based on the kernels of the contraction maps $\left\{  \left.
L_{N}^{P}\right\vert 0\leq p\leq N\right\}  ,$ namely
\begin{equation}
\mathcal{K}_{Np}=\ker L_{N}^{p-1}\ominus\ker L_{N}^{p};\quad1\leq p\leq N
\end{equation}
and
\begin{equation}
\mathcal{B}_{1}\left(  \mathcal{H}^{N}\right)  =%
%TCIMACRO{\tbigoplus \limits_{1\leq p\leq N}}%
%BeginExpansion
{\textstyle\bigoplus\limits_{1\leq p\leq N}}
%EndExpansion
\mathcal{K}_{Np}\oplus\mathcal{K}^{c}%
\end{equation}
where $\mathcal{K}^{c}$ is a compliment of $%
%TCIMACRO{\tbigoplus \limits_{1\leq p\leq N}}%
%BeginExpansion
{\textstyle\bigoplus\limits_{1\leq p\leq N}}
%EndExpansion
\mathcal{K}_{Np}$ in $\mathcal{B}_{1}\left(  \mathcal{H}^{N}\right)  $.

\subsection{Examples}

Consider the bases of the operator space $\mathcal{B}_{1}\left(
\mathcal{H}^{1}\right)  $ given by%
\begin{equation}
\left\{  \left\{  \left.  \left\vert \varphi_{j}\right\rangle \left\langle
\varphi_{j}\right\vert -\left\vert \varphi_{j+1}\right\rangle \left\langle
\varphi_{j+1}\right\vert \right\vert 1\leq j\leq\infty\right\}  ,\left\{
\left.  \left\vert \varphi_{j}\right\rangle \left\langle \varphi
_{k}\right\vert \right\vert 1\leq j\neq k\leq\infty\right\}  ,\left\vert
\varphi_{1}\right\rangle \left\langle \varphi_{1}\right\vert \right\}
\end{equation}
that is adapted to the decomposition
\begin{equation}
\mathcal{B}_{1}\left(  \mathcal{H}^{1}\right)  =\mathcal{K}_{11od}%
\oplus\mathcal{K}_{11d}\oplus\mathcal{S}_{1d} \label{graddec}%
\end{equation}
where
\begin{align}
\mathcal{K}_{11od}  &  =\overline{\operatorname{linspan}}\left\{  \left.
\left\vert \varphi_{j}\right\rangle \left\langle \varphi_{k}\right\vert
\right\vert 1\leq j\neq k\leq\infty\right\} \nonumber\\
\mathcal{K}_{11d}  &  =\overline{\operatorname{linspan}}\left\{  \left.
\left\vert \varphi_{j}\right\rangle \left\langle \varphi_{j}\right\vert
-\left\vert \varphi_{j+1}\right\rangle \left\langle \varphi_{j+1}\right\vert
\right\vert 1\leq j\leq\infty\right\} \nonumber\\
\mathcal{S}_{1d}  &  =\left\{  \left.  \alpha\left\vert \varphi_{1}%
\right\rangle \left\langle \varphi_{1}\right\vert \right\vert \alpha
\in\mathbb{C}\right\}
\end{align}
Let us denote these basis elements by
\begin{align}
\left\vert \varphi_{jk}\right)   &  =\left\vert \varphi_{j}\right\rangle
\left\langle \varphi_{k}\right\vert ;\quad1\leq j\neq k\leq\infty\nonumber\\
\left\vert \varphi_{j+1}\right)   &  =\left\vert \varphi_{j}\right\rangle
\left\langle \varphi_{j}\right\vert -\left\vert \varphi_{j+1}\right\rangle
\left\langle \varphi_{j+1}\right\vert ;\quad1\leq j\leq\infty\nonumber\\
\left\vert \varphi_{1}\right)   &  =\left\vert \varphi_{1}\right\rangle
\left\langle \varphi_{1}\right\vert
\end{align}
We can note that
\begin{equation}
\left\vert \varphi_{1}\right\rangle \left\langle \varphi_{1}\right\vert
-\left(  \left\vert \varphi_{1}\right\rangle \left\langle \varphi
_{1}\right\vert -\left\vert \varphi_{2}\right\rangle \left\langle \varphi
_{2}\right\vert \right)  =\left\vert \varphi_{2}\right\rangle \left\langle
\varphi_{2}\right\vert
\end{equation}
and
\begin{equation}
\left(  \left\vert \varphi_{1}\right\rangle \left\langle \varphi
_{1}\right\vert -\left\vert \varphi_{2}\right\rangle \left\langle \varphi
_{2}\right\vert \right)  +\left(  \left\vert \varphi_{2}\right\rangle
\left\langle \varphi_{2}\right\vert -\left\vert \varphi_{3}\right\rangle
\left\langle \varphi_{3}\right\vert \right)  =\left(  \left\vert \varphi
_{1}\right\rangle \left\langle \varphi_{1}\right\vert -\left\vert \varphi
_{3}\right\rangle \left\langle \varphi_{3}\right\vert \right)
\end{equation}
so that
\begin{equation}
\left\vert \varphi_{1}\right\rangle \left\langle \varphi_{1}\right\vert
-\left(  \left\vert \varphi_{1}\right\rangle \left\langle \varphi
_{1}\right\vert -\left\vert \varphi_{2}\right\rangle \left\langle \varphi
_{2}\right\vert \right)  -\left(  \left\vert \varphi_{2}\right\rangle
\left\langle \varphi_{2}\right\vert -\left\vert \varphi_{3}\right\rangle
\left\langle \varphi_{3}\right\vert \right)  =\left\vert \varphi
_{3}\right\rangle \left\langle \varphi_{3}\right\vert
\end{equation}
etc., hence generating all diagonal basis elements $\left\{  \left.
\left\vert \varphi_{j}\right\rangle \left\langle \varphi_{j}\right\vert
\right\vert ;2\leq j\leq\infty\right\}  $ as%
\begin{equation}
\left\vert \varphi_{j}\right\rangle \left\langle \varphi_{j}\right\vert
=\left\vert \varphi_{1}\right)  -\sum_{2\leq m\leq j}\left\vert \varphi
_{j}\right)  ;\quad1\leq j\leq\infty
\end{equation}
Any element of $X\in\mathcal{B}_{1}\left(  \mathcal{H}^{1}\right)  $ can be
expanded%
\begin{align}
A  &  =\sum_{1\leq j,k\leq\infty}A_{jk}\left\vert \varphi_{j}\right\rangle
\left\langle \varphi_{k}\right\vert \\
&  =\sum_{1\leq j\neq k\leq\infty}X_{jk}\left\vert \varphi_{jk}\right)
+\sum_{2\leq j\leq\infty}X_{\mathcal{K}j}\left\vert \varphi_{j}\right)
+X_{1}\left\vert \varphi_{1}\right)
\end{align}
where
\begin{align}
X_{1}  &  =Tr\left\{  A\right\} \nonumber\\
X_{jk}  &  =A_{jk}\nonumber\\
X_{j}  &  =Tr\left\{  A\right\}  -\sum_{1\leq m\leq j-1}A_{mm};\quad2\leq
j\leq\infty
\end{align}
The choice of $\left\vert \varphi_{1}\right)  =\left\vert \varphi
_{1}\right\rangle \left\langle \varphi_{1}\right\vert $ is quite arbitray, one
can just as well use $\left\vert \varphi_{1}\right)  =\left\vert \varphi
_{j}\right\rangle \left\langle \varphi_{j}\right\vert $ for any other $j,$
which just produces a change of basis for $\mathcal{B}_{1}\left(
\mathcal{H}^{N}\right)  $ and hence a change of expansion coefficients while
still produces a decomposition that is isomorphic to eq. $\left(
\ref{graddec}\right)  $.

One can use the basis $\left\{  \left\{  \left.  \left\vert \varphi
_{j}\right)  \right\vert 1\leq j\leq\infty\right\}  ,\left\{  \left.
\left\vert \varphi_{jk}\right)  \right\vert 1\leq j\neq k\leq\infty\right\}
\right\}  $ to construct basis of $\mathcal{B}_{1}\left(  \mathcal{H}%
^{N}\right)  $ for $2\leq N\leq\infty$ in the following manner:

$\mathcal{B}_{1}\left(  \mathcal{H}^{2}\right)  $ can be decomposed as
\begin{align}
\mathcal{B}_{1}\left(  \mathcal{H}^{2}\right)   &  =\left(  \ker L_{2}%
^{1}\ominus\ker L_{2}^{2}\right)  \oplus\left(  \ker L_{2}^{0}\ominus\ker
L_{2}^{1}\right)  \oplus\left(  \ker L_{2}^{0}\right)  ^{c}\nonumber\\
&  =\ker L_{2}^{0}\oplus\left(  \ker L_{2}^{0}\right)  ^{c}%
\end{align}
as $\ker L_{2}^{2}=\phi$, in terms of the notation introduced this is
\begin{equation}
\mathcal{B}_{1}\left(  \mathcal{H}^{2}\right)  =\mathcal{K}_{22}%
\oplus\mathcal{K}_{21}\oplus\mathcal{K}_{20}%
\end{equation}
where $\mathcal{K}_{20}=$ $\mathcal{K}^{c}$ in $\mathcal{B}_{1}\left(
\mathcal{H}^{2}\right)  .$ It can be shown that
\begin{equation}
\mathcal{K}_{21}\oplus\mathcal{K}_{20}\cong\mathcal{B}_{1}\left(
\mathcal{H}^{1}\right)
\end{equation}
This decomposition can be analysed in terms of lie groups as
\begin{equation}
\mathcal{B}_{1}\left(  \mathcal{H}^{2}\right)  =sl\left(  \mathcal{H}%
^{2}\right)  _{2}\oplus sl\left(  \mathcal{H}^{2}\right)  _{2}^{c}%
\end{equation}
as
\begin{align}
sl\left(  \mathcal{H}^{2}\right)  _{2} &  =\mathcal{K}_{22}\oplus
\mathcal{K}_{21}\nonumber\\
sl\left(  \mathcal{H}^{2}\right)  _{2}^{c} &  =\mathcal{K}_{20}%
\end{align}
and
\begin{equation}
\mathcal{B}_{1}\left(  \mathcal{H}^{2}\right)  =\left(  sl\left(
\mathcal{H}^{2}\right)  _{2}\ominus sl\left(  \mathcal{H}^{1}\right)
_{2}\right)  \oplus sl\left(  \mathcal{H}^{1}\right)  _{2}\oplus sl\left(
\mathcal{H}^{2}\right)  _{2}^{c}%
\end{equation}
where $sl\left(  \mathcal{H}^{2}\right)  _{2}$ and $sl\left(  \mathcal{H}%
^{1}\right)  _{2}$ are representations of $sl\left(  \mathcal{H}^{2}\right)  $
and $sl\left(  \mathcal{H}^{1}\right)  $ respectively in $\mathcal{H}^{2}.$ By
exponentiating these subspaces one obtains representations of the lie groups
\begin{align}
SL\left(  \mathcal{H}^{2}\right)  _{2} &  =\exp\left(  sl\left(
\mathcal{H}^{2}\right)  _{2}\right)  \nonumber\\
SL\left(  \mathcal{H}^{1}\right)  _{2} &  =\exp\left(  sl\left(
\mathcal{H}^{1}\right)  _{2}\right)
\end{align}
This decomposition also has the property that it carries an adjoint
representation of $SU\left(  \mathcal{H}^{1}\right)  $ i.e.%
\begin{align}
\operatorname{Adj}\left(  SU\left(  \mathcal{H}^{1}\right)  _{2}\right)   &
:\mathcal{K}_{22}\rightarrow\mathcal{K}_{22}\nonumber\\
\operatorname{Adj}\left(  SU\left(  \mathcal{H}^{1}\right)  _{2}\right)   &
:\mathcal{K}_{21}\rightarrow\mathcal{K}_{21}%
\end{align}
which can be deduced from the intertwining property
\begin{equation}
L_{N}^{P}\circ\operatorname{Adj}\left(  SU\left(  \mathcal{H}^{1}\right)
_{p}\right)  =\operatorname{Adj}\left(  SU\left(  \mathcal{H}^{1}\right)
_{p}\right)  \circ L_{N}^{P}%
\end{equation}
which gives that
\begin{equation}
\operatorname{Adj}\left(  SU\left(  \mathcal{H}^{1}\right)  _{p}\right)  :\ker
L_{N}^{P}\rightarrow\ker L_{N}^{P}%
\end{equation}
so in particular
\begin{equation}
L_{2}^{1}\circ\operatorname{Adj}\left(  SU\left(  \mathcal{H}^{1}\right)
_{1}\right)  =\operatorname{Adj}\left(  SU\left(  \mathcal{H}^{1}\right)
_{1}\right)  \circ L_{2}^{1}%
\end{equation}
The subspaces $\mathcal{K}_{22}$ and $\mathcal{K}_{21}$ can be decomposed into
diagonal and non diagonal parts wrt the basis $\left\{  \left\vert
\varphi_{j_{1}}\varphi_{j_{2}}\right\rangle \right\}  $ of $\mathcal{H}^{2}$
i.e.
\begin{align}
\mathcal{K}_{22} &  =\mathcal{K}_{22od}\oplus\mathcal{K}_{22d}\nonumber\\
\mathcal{K}_{21} &  =\mathcal{K}_{21od}\oplus\mathcal{K}_{21d}%
\end{align}
so that
\begin{equation}
\mathcal{B}_{1}\left(  \mathcal{H}^{2}\right)  =\mathcal{K}_{22od}%
\oplus\mathcal{K}_{22d}\oplus\mathcal{K}_{21od}\oplus\mathcal{K}_{21d}%
\oplus\mathcal{K}_{20}%
\end{equation}
We can consider as an example a space $\mathcal{H}^{1}$ that has dimension 4.
Then%
\begin{equation}
\left(  \ker L_{2}^{1}\ominus\ker L_{2}^{2}\right)  \oplus\left(  \ker
L_{2}^{0}\ominus\ker L_{2}^{1}\right)  \oplus\left(  \ker L_{2}^{0}\right)
^{c}\equiv\mathcal{K}_{22}\oplus\mathcal{K}_{21}\oplus\mathcal{K}_{20}%
\end{equation}
and
\begin{align}
&  \dim\left(  \mathcal{K}_{21}\oplus\mathcal{K}_{20}\right)  =\dim\left(
\mathcal{B}_{1}\left(  \mathcal{H}^{1}\right)  \right)  =16\nonumber\\
&  \dim\left(  \mathcal{K}_{20}\right)  =1\nonumber\\
&  \dim\left(  \mathcal{K}_{22}\right)  =\dim\left(  \mathcal{B}_{1}\left(
\mathcal{H}^{2}\right)  \right)  -\dim\left(  \mathcal{B}_{1}\left(
\mathcal{H}^{1}\right)  \right)  =36-16=20\nonumber\\
&  \dim\left(  \mathcal{K}_{22od}\right)  =18\nonumber\\
&  \dim\left(  \mathcal{K}_{22d}\right)  =2\nonumber\\
&  \dim\left(  \mathcal{K}_{21od}\right)  =12\nonumber\\
&  \dim\left(  \mathcal{K}_{21d}\right)  =3
\end{align}


The space $\mathcal{K}_{22od}$ is generated by the $18$ linear independent
elements (i.e. is of dimension $18)$
\begin{equation}
\mathfrak{b}_{0}\left(  \mathbf{\varphi}\right)  =\left\{  \left\vert
\varphi_{1}\varphi_{2}\right\rangle \left\langle \varphi_{3}\varphi
_{4}\right\vert ,\left\vert \varphi_{1}\varphi_{3}\right\rangle \left\langle
\varphi_{2}\varphi_{4}\right\vert ,\left\vert \varphi_{1}\varphi
_{4}\right\rangle \left\langle \varphi_{2}\varphi_{3}\right\vert ,\left\vert
\varphi_{2}\varphi_{3}\right\rangle \left\langle \varphi_{1}\varphi
_{4}\right\vert ,\left\vert \varphi_{2}\varphi_{4}\right\rangle \left\langle
\varphi_{1}\varphi_{3}\right\vert ,\left\vert \varphi_{3}\varphi
_{4}\right\rangle \left\langle \varphi_{1}\varphi_{2}\right\vert \right\}
\end{equation}
\begin{align}
& \mathbf{\lceil}\text{In general }\left\{  \left\vert \varphi_{j_{1}%
k_{1}j_{2}k_{2}}^{0}\right)  =\left\vert \varphi_{j_{1}k_{1}}\right)
\wedge\left\vert \varphi_{j_{2}k_{2}}\right)  =\left\vert \varphi_{j_{1}%
}\right\rangle \left\langle \varphi_{k_{1}}\right\vert \wedge\left\vert
\varphi_{j_{2}}\right\rangle \left\langle \varphi_{k_{2}}\right\vert
=\left\vert \varphi_{j_{1}}\varphi_{j_{2}}\right\rangle \left\langle
\varphi_{k_{1}}\varphi_{k_{2}}\right\vert \right.  \nonumber\\
& \qquad\qquad\qquad\qquad\qquad\qquad\qquad\qquad\qquad\qquad\left.
j_{1}\neq k_{1},k_{2};\quad j_{2}\neq k_{1},k_{2};\quad j_{1}<j_{2}%
;k_{1}<k_{2}\right\}  \rfloor
\end{align}


and
\begin{align}
&  \left\{  \left(  \left\vert \varphi_{1}\right\rangle \left\langle
\varphi_{1}\right\vert -\left\vert \varphi_{2}\right\rangle \left\langle
\varphi_{2}\right\vert \right)  \wedge\left\vert \varphi_{3}\right\rangle
\left\langle \varphi_{4}\right\vert ,\left(  \left\vert \varphi_{1}%
\right\rangle \left\langle \varphi_{1}\right\vert -\left\vert \varphi
_{2}\right\rangle \left\langle \varphi_{2}\right\vert \right)  \wedge
\left\vert \varphi_{4}\right\rangle \left\langle \varphi_{3}\right\vert
\right.  \nonumber\\
&  \,\,\,\left(  \left\vert \varphi_{1}\right\rangle \left\langle \varphi
_{1}\right\vert -\left\vert \varphi_{3}\right\rangle \left\langle \varphi
_{3}\right\vert \right)  \wedge\left\vert \varphi_{2}\right\rangle
\left\langle \varphi_{4}\right\vert ,\left(  \left\vert \varphi_{1}%
\right\rangle \left\langle \varphi_{1}\right\vert -\left\vert \varphi
_{3}\right\rangle \left\langle \varphi_{3}\right\vert \right)  \wedge
\left\vert \varphi_{2}\right\rangle \left\langle \varphi_{4}\right\vert
\nonumber\\
&  \,\,\,\left(  \left\vert \varphi_{1}\right\rangle \left\langle \varphi
_{1}\right\vert -\left\vert \varphi_{4}\right\rangle \left\langle \varphi
_{4}\right\vert \right)  \wedge\left\vert \varphi_{2}\right\rangle
\left\langle \varphi_{3}\right\vert ,\left(  \left\vert \varphi_{1}%
\right\rangle \left\langle \varphi_{1}\right\vert -\left\vert \varphi
_{4}\right\rangle \left\langle \varphi_{4}\right\vert \right)  \wedge
\left\vert \varphi_{3}\right\rangle \left\langle \varphi_{2}\right\vert
\nonumber\\
&  \,\,\,\left(  \left\vert \varphi_{2}\right\rangle \left\langle \varphi
_{2}\right\vert -\left\vert \varphi_{3}\right\rangle \left\langle \varphi
_{3}\right\vert \right)  \wedge\left\vert \varphi_{1}\right\rangle
\left\langle \varphi_{4}\right\vert ,\left(  \left\vert \varphi_{2}%
\right\rangle \left\langle \varphi_{2}\right\vert -\left\vert \varphi
_{3}\right\rangle \left\langle \varphi_{3}\right\vert \right)  \wedge
\left\vert \varphi_{4}\right\rangle \left\langle \varphi_{1}\right\vert
\nonumber\\
&  \,\,\,\left(  \left\vert \varphi_{2}\right\rangle \left\langle \varphi
_{2}\right\vert -\left\vert \varphi_{4}\right\rangle \left\langle \varphi
_{4}\right\vert \right)  \wedge\left\vert \varphi_{1}\right\rangle
\left\langle \varphi_{3}\right\vert ,\left(  \left\vert \varphi_{2}%
\right\rangle \left\langle \varphi_{2}\right\vert -\left\vert \varphi
_{4}\right\rangle \left\langle \varphi_{4}\right\vert \right)  \wedge
\left\vert \varphi_{3}\right\rangle \left\langle \varphi_{1}\right\vert
\nonumber\\
&  \,\left.  \left(  \left\vert \varphi_{3}\right\rangle \left\langle
\varphi_{3}\right\vert -\left\vert \varphi_{4}\right\rangle \left\langle
\varphi_{4}\right\vert \right)  \wedge\left\vert \varphi_{1}\right\rangle
\left\langle \varphi_{2}\right\vert ,\left(  \left\vert \varphi_{3}%
\right\rangle \left\langle \varphi_{3}\right\vert -\left\vert \varphi
_{4}\right\rangle \left\langle \varphi_{4}\right\vert \right)  \wedge
\left\vert \varphi_{2}\right\rangle \left\langle \varphi_{1}\right\vert
\right\}  \subset\mathfrak{b}_{1}\left(  \mathbf{\varphi}\right)
\end{align}
The space $\mathcal{K}_{22d}$ is generated by
\begin{subequations}
\begin{align}
&  \left(  \left\vert \varphi_{1}\right\rangle \left\langle \varphi
_{1}\right\vert -\left\vert \varphi_{2}\right\rangle \left\langle \varphi
_{2}\right\vert \right)  \wedge\left(  \left\vert \varphi_{3}\right\rangle
\left\langle \varphi_{3}\right\vert -\left\vert \varphi_{4}\right\rangle
\left\langle \varphi_{4}\right\vert \right)  \nonumber\\
&  =\left\vert \varphi_{1}\varphi_{3}\right\rangle \left\langle \varphi
_{1}\varphi_{3}\right\vert -\left\vert \varphi_{2}\varphi_{3}\right\rangle
\left\langle \varphi_{2}\varphi_{3}\right\vert -\left\vert \varphi_{1}%
\varphi_{4}\right\rangle \left\langle \varphi_{1}\varphi_{4}\right\vert
+\left\vert \varphi_{2}\varphi_{4}\right\rangle \left\langle \varphi
_{2}\varphi_{4}\right\vert \label{b1}\\
&  \left(  \left\vert \varphi_{1}\right\rangle \left\langle \varphi
_{1}\right\vert -\left\vert \varphi_{3}\right\rangle \left\langle \varphi
_{3}\right\vert \right)  \wedge\left(  \left\vert \varphi_{2}\right\rangle
\left\langle \varphi_{2}\right\vert -\left\vert \varphi_{4}\right\rangle
\left\langle \varphi_{4}\right\vert \right)  \nonumber\\
&  =\left\vert \varphi_{1}\varphi_{2}\right\rangle \left\langle \varphi
_{1}\varphi_{2}\right\vert -\left\vert \varphi_{2}\varphi_{3}\right\rangle
\left\langle \varphi_{2}\varphi_{3}\right\vert -\left\vert \varphi_{1}%
\varphi_{4}\right\rangle \left\langle \varphi_{1}\varphi_{4}\right\vert
+\left\vert \varphi_{3}\varphi_{4}\right\rangle \left\langle \varphi
_{3}\varphi_{4}\right\vert \label{b2}\\
&  \left(  \left\vert \varphi_{1}\right\rangle \left\langle \varphi
_{1}\right\vert -\left\vert \varphi_{4}\right\rangle \left\langle \varphi
_{4}\right\vert \right)  \wedge\left(  \left\vert \varphi_{2}\right\rangle
\left\langle \varphi_{2}\right\vert -\left\vert \varphi_{3}\right\rangle
\left\langle \varphi_{3}\right\vert \right)  \nonumber\\
&  =\left\vert \varphi_{1}\varphi_{2}\right\rangle \left\langle \varphi
_{1}\varphi_{2}\right\vert -\left\vert \varphi_{2}\varphi_{4}\right\rangle
\left\langle \varphi_{2}\varphi_{4}\right\vert -\left\vert \varphi_{1}%
\varphi_{3}\right\rangle \left\langle \varphi_{1}\varphi_{3}\right\vert
+\left\vert \varphi_{3}\varphi_{4}\right\rangle \left\langle \varphi
_{3}\varphi_{4}\right\vert \label{b3}%
\end{align}
However basis element eq. $\left(  \ref{b1}\right)  -$ basis element eq.
$\left(  \ref{b2}\right)  =$ basis element eq. $\left(  \ref{b3}\right)  $, so
we can choose basis element eq. $\left(  \ref{b1}\right)  $ and eq. $\left(
\ref{b2}\right)  $
\end{subequations}
\begin{align}
&  \left\{  \left(  \left\vert \varphi_{1}\right\rangle \left\langle
\varphi_{1}\right\vert -\left\vert \varphi_{2}\right\rangle \left\langle
\varphi_{2}\right\vert \right)  \wedge\left(  \left\vert \varphi
_{3}\right\rangle \left\langle \varphi_{3}\right\vert -\left\vert \varphi
_{4}\right\rangle \left\langle \varphi_{4}\right\vert \right)  ,\right.
\nonumber\\
&  \left.  \left(  \left\vert \varphi_{1}\right\rangle \left\langle
\varphi_{1}\right\vert -\left\vert \varphi_{3}\right\rangle \left\langle
\varphi_{3}\right\vert \right)  \wedge\left(  \left\vert \varphi
_{2}\right\rangle \left\langle \varphi_{2}\right\vert -\left\vert \varphi
_{4}\right\rangle \left\langle \varphi_{4}\right\vert \right)  \right\}
\subset\mathfrak{b}_{1}\left(  \mathbf{\varphi}\right)
\end{align}%
\begin{align}
& \mathbf{\lceil}\text{In general }\left\{  \left(  \left\vert \varphi_{j_{1}%
}\right\rangle \left\langle \varphi_{j_{1}}\right\vert -\left\vert
\varphi_{j_{2}}\right\rangle \left\langle \varphi_{j_{2}}\right\vert \right)
\wedge\left(  \left\vert \varphi_{j_{3}}\right\rangle \left\langle
\varphi_{j_{3}}\right\vert -\left\vert \varphi_{j_{4}}\right\rangle
\left\langle \varphi_{j_{4}}\right\vert \right)  \right.  =\nonumber\\
& \,\left\vert \varphi_{j_{1}}\varphi_{j_{3}}\right\rangle \left\langle
\varphi_{j_{1}}\varphi_{j_{3}}\right\vert -\left\vert \varphi_{j_{1}}%
\varphi_{j_{4}}\right\rangle \left\langle \varphi_{j_{1}}\varphi_{j_{4}%
}\right\vert -\left\vert \varphi_{j_{2}}\varphi_{j_{3}}\right\rangle
\left\langle \varphi_{j_{2}}\varphi_{j_{3}}\right\vert +\left\vert
\varphi_{j_{2}}\varphi_{j_{4}}\right\rangle \left\langle \varphi_{j_{2}%
}\varphi_{j_{4}}\right\vert \qquad\left.  j_{1}\neq j_{2}\neq j_{3}\neq
j_{4}\right\}  \rfloor
\end{align}
to generate $\mathcal{K}_{22d}$ which is evidently of dimension $2.$We can
check the result of the contraction map $L_{2}^{1}$ on these elements by
evaluating generic sample elements
\[
Tr\left\{  \left(  \left\vert \varphi_{1}\varphi_{3}\right\rangle \left\langle
\varphi_{1}\varphi_{4}\right\vert -\left\vert \varphi_{2}\varphi
_{3}\right\rangle \left\langle \varphi_{2}\varphi_{4}\right\vert \right)
a_{j}^{\dagger}a_{k}\right\}  =0\quad\forall j,k
\]
and
\begin{equation}
Tr\left\{  \left\vert \varphi_{1}\varphi_{2}\right\rangle \left\langle
\varphi_{3}\varphi_{4}\right\vert a_{j}^{\dagger}a_{k}\right\}  =0\quad\forall
j,k
\end{equation}
and
\[
Tr\left\{  \left(  \left\vert \varphi_{1}\varphi_{3}\right\rangle \left\langle
\varphi_{1}\varphi_{3}\right\vert -\left\vert \varphi_{1}\varphi
_{4}\right\rangle \left\langle \varphi_{1}\varphi_{4}\right\vert -\left\vert
\varphi_{2}\varphi_{3}\right\rangle \left\langle \varphi_{2}\varphi
_{3}\right\vert +\left\vert \varphi_{2}\varphi_{4}\right\rangle \left\langle
\varphi_{2}\varphi_{4}\right\vert \right)  a_{j}^{\dagger}a_{k}\right\}
=0\quad\forall j,k
\]
Obviously such elements $\notin\ker L_{2}^{2}$ $=\phi$ as $\left\{  Tr\left\{
Ba_{j_{1}}^{\dagger}a_{j_{2}}^{\dagger}a_{k_{2}}a_{k_{1}}\right\}  \right\}  $
are not all zero.

The space $\mathcal{K}_{21od}$ of dimension $12$ is generated by
\begin{align}
&  \left\{  \left(  \left\vert \varphi_{1}\right\rangle \left\langle
\varphi_{1}\right\vert +\left\vert \varphi_{2}\right\rangle \left\langle
\varphi_{2}\right\vert \right)  \wedge\left\vert \varphi_{3}\right\rangle
\left\langle \varphi_{4}\right\vert ,\left(  \left\vert \varphi_{1}%
\right\rangle \left\langle \varphi_{1}\right\vert +\left\vert \varphi
_{2}\right\rangle \left\langle \varphi_{2}\right\vert \right)  \wedge
\left\vert \varphi_{4}\right\rangle \left\langle \varphi_{3}\right\vert
\right.  \nonumber\\
&  \left(  \,\,\,\left\vert \varphi_{1}\right\rangle \left\langle \varphi
_{1}\right\vert +\left\vert \varphi_{3}\right\rangle \left\langle \varphi
_{3}\right\vert \right)  \wedge\left\vert \varphi_{2}\right\rangle
\left\langle \varphi_{4}\right\vert ,\left(  \left\vert \varphi_{1}%
\right\rangle \left\langle \varphi_{1}\right\vert +\left\vert \varphi
_{3}\right\rangle \left\langle \varphi_{3}\right\vert \right)  \wedge
\left\vert \varphi_{2}\right\rangle \left\langle \varphi_{4}\right\vert
\nonumber\\
&  \left(  \,\,\,\left\vert \varphi_{1}\right\rangle \left\langle \varphi
_{1}\right\vert +\left\vert \varphi_{4}\right\rangle \left\langle \varphi
_{4}\right\vert \right)  \wedge\left\vert \varphi_{2}\right\rangle
\left\langle \varphi_{3}\right\vert ,\left(  \left\vert \varphi_{1}%
\right\rangle \left\langle \varphi_{1}\right\vert +\left\vert \varphi
_{4}\right\rangle \left\langle \varphi_{4}\right\vert \right)  \wedge
\left\vert \varphi_{3}\right\rangle \left\langle \varphi_{2}\right\vert
\nonumber\\
&  \left(  \,\,\,\left\vert \varphi_{2}\right\rangle \left\langle \varphi
_{2}\right\vert +\left\vert \varphi_{3}\right\rangle \left\langle \varphi
_{3}\right\vert \right)  \wedge\left\vert \varphi_{1}\right\rangle
\left\langle \varphi_{4}\right\vert ,\left(  \left\vert \varphi_{2}%
\right\rangle \left\langle \varphi_{2}\right\vert +\left\vert \varphi
_{3}\right\rangle \left\langle \varphi_{3}\right\vert \right)  \wedge
\left\vert \varphi_{4}\right\rangle \left\langle \varphi_{1}\right\vert
\nonumber\\
&  \left(  \,\,\,\left\vert \varphi_{2}\right\rangle \left\langle \varphi
_{2}\right\vert +\left\vert \varphi_{4}\right\rangle \left\langle \varphi
_{4}\right\vert \right)  \wedge\left\vert \varphi_{1}\right\rangle
\left\langle \varphi_{3}\right\vert ,\left(  \left\vert \varphi_{2}%
\right\rangle \left\langle \varphi_{2}\right\vert +\left\vert \varphi
_{4}\right\rangle \left\langle \varphi_{4}\right\vert \right)  \wedge
\left\vert \varphi_{3}\right\rangle \left\langle \varphi_{1}\right\vert
\nonumber\\
&  \left.  \left(  \,\left\vert \varphi_{3}\right\rangle \left\langle
\varphi_{3}\right\vert +\left\vert \varphi_{4}\right\rangle \left\langle
\varphi_{4}\right\vert \right)  \wedge\left\vert \varphi_{1}\right\rangle
\left\langle \varphi_{2}\right\vert ,\left(  \left\vert \varphi_{3}%
\right\rangle \left\langle \varphi_{3}\right\vert +\left\vert \varphi
_{4}\right\rangle \left\langle \varphi_{4}\right\vert \right)  \wedge
\left\vert \varphi_{2}\right\rangle \left\langle \varphi_{1}\right\vert
\right\}  \subset\mathfrak{b}_{1}\left(  \mathbf{\varphi}\right)
\end{align}%
\begin{align}
& \mathbf{\lceil}\text{In general }\left\{  \left(  \left\vert \varphi_{j_{1}%
}\right\rangle \left\langle \varphi_{j_{1}}\right\vert +\left\vert
\varphi_{j_{2}}\right\rangle \left\langle \varphi_{j_{2}}\right\vert \right)
\wedge\left\vert \varphi_{j_{3}}\right\rangle \left\langle \varphi_{j_{4}%
}\right\vert \right.  =\nonumber\\
& \,\left\vert \varphi_{j_{1}}\varphi_{j_{3}}\right\rangle \left\langle
\varphi_{j_{1}}\varphi_{j_{4}}\right\vert +\left\vert \varphi_{j_{2}}%
\varphi_{j_{3}}\right\rangle \left\langle \varphi_{j_{2}}\varphi_{j_{4}%
}\right\vert \qquad\left.  \left(  j_{1}<j_{2}\right)  \neq\left(  j_{3}\neq
j_{4}\right)  \right\}  \rfloor
\end{align}
The space $\mathcal{K}_{21d}$ must be of dimension $3$ as $\dim\left(
\mathcal{K}_{22d}\right)  =2$ and $\dim\left(  \mathcal{K}_{20}\right)  =1$
and has a basis given by
\begin{align}
&  \left\{  \left\vert \varphi_{1}\right\rangle \left\langle \varphi
_{1}\right\vert \wedge\left(  \left\vert \varphi_{2}\right\rangle \left\langle
\varphi_{2}\right\vert -\left\vert \varphi_{3}\right\rangle \left\langle
\varphi_{3}\right\vert \right)  ,\right.  \nonumber\\
&  \,\,\,\left\vert \varphi_{1}\right\rangle \left\langle \varphi
_{1}\right\vert \wedge\left(  \left\vert \varphi_{2}\right\rangle \left\langle
\varphi_{2}\right\vert -\left\vert \varphi_{4}\right\rangle \left\langle
\varphi_{4}\right\vert \right)  \nonumber\\
&  \left.  \,\left\vert \varphi_{1}\right\rangle \left\langle \varphi
_{1}\right\vert \wedge\left(  \left\vert \varphi_{3}\right\rangle \left\langle
\varphi_{3}\right\vert -\left\vert \varphi_{4}\right\rangle \left\langle
\varphi_{4}\right\vert \right)  \right\}  \subset\mathfrak{b}_{1}\left(
\mathbf{\varphi}\right)
\end{align}%
\begin{equation}
\mathbf{\lceil}\text{In general }\left\{  \left\vert \varphi_{1}\right\rangle
\left\langle \varphi_{1}\right\vert \wedge\left(  \left\vert \varphi_{j_{1}%
}\right\rangle \left\langle \varphi_{j_{1}}\right\vert -\left\vert
\varphi_{j_{2}}\right\rangle \left\langle \varphi_{j_{2}}\right\vert \right)
\right.  =\left\vert \varphi_{1}\varphi_{j_{1}}\right\rangle \left\langle
\varphi_{1}\varphi_{j_{1}}\right\vert -\left\vert \varphi_{1}\varphi_{j_{2}%
}\right\rangle \left\langle \varphi_{1}\varphi_{j_{2}}\right\vert
\qquad\left.  j_{1}\neq j_{2}\neq1\right\}  \rfloor
\end{equation}
We can check the result of the contraction map $L_{2}^{0}$ on these elements
by evaluating
\begin{equation}
Tr\left\{  B\right\}
\end{equation}
which gives for example
\begin{align*}
Tr\left\{  \left\vert \varphi_{1}\varphi_{2}\right\rangle \left\langle
\varphi_{1}\varphi_{3}\right\vert \right\}   &  =0\\
Tr\left\{  \left\vert \varphi_{1}\varphi_{2}\right\rangle \left\langle
\varphi_{1}\varphi_{2}\right\vert -\left\vert \varphi_{1}\varphi
_{3}\right\rangle \left\langle \varphi_{1}\varphi_{3}\right\vert \right\}   &
=0
\end{align*}
and the results of contraction map $L_{2}^{1}$ by
\begin{equation}
Tr\left\{  Ba_{j}^{\dagger}a_{k}\right\}
\end{equation}
which gives for example%
\begin{align*}
Tr\left\{  \left\vert \varphi_{1}\varphi_{2}\right\rangle \left\langle
\varphi_{1}\varphi_{3}\right\vert a_{j}^{\dagger}a_{k}\right\}   &
=0\quad\not \forall\,j,k\\
Tr\left\{  \left(  \left\vert \varphi_{1}\varphi_{2}\right\rangle \left\langle
\varphi_{1}\varphi_{2}\right\vert -\left\vert \varphi_{1}\varphi
_{3}\right\rangle \left\langle \varphi_{1}\varphi_{3}\right\vert \right)
a_{j}^{\dagger}a_{k}\right\}   &  =0\quad\not \forall\,j,k
\end{align*}


The space $\mathcal{S}_{2d}$ of dimension $1$ is spanned by
\begin{equation}
\left\vert \varphi_{1}\varphi_{2}\right\rangle \left\langle \varphi_{1}%
\varphi_{2}\right\vert =\left\vert \varphi_{1}\right\rangle \left\langle
\varphi_{1}\right\vert \wedge\left\vert \varphi_{2}\right\rangle \left\langle
\varphi_{2}\right\vert
\end{equation}


$\mathcal{B}_{1}\left(  \mathcal{H}^{3}\right)  $ can be decomposed as
\begin{align}
\mathcal{B}_{1}\left(  \mathcal{H}^{3}\right)   &  =\left(  \ker L_{3}%
^{2}\ominus\ker L_{3}^{3}\right)  \oplus\left(  \ker L_{3}^{1}\ominus\ker
L_{3}^{2}\right)  \oplus\left(  \ker L_{3}^{0}\ominus\ker L_{3}^{1}\right)
\oplus\left(  \ker L_{3}^{0}\right)  ^{c}\nonumber\\
&  =\ker L_{3}^{0}\oplus\left(  \ker L_{2}^{0}\right)  ^{c}\nonumber\\
&  \equiv\mathcal{K}_{33}\oplus\mathcal{K}_{32}\oplus\mathcal{K}_{31}%
\oplus\mathcal{K}_{30}%
\end{align}
\bigskip We have the following isomorphims%
\begin{align}
\mathcal{K}_{31}\oplus\mathcal{K}_{30}  & \cong\mathcal{B}_{1}\left(
\mathcal{H}^{1}\right)  \nonumber\\
\mathcal{K}_{32}\oplus\mathcal{K}_{31}\oplus\mathcal{K}_{30}  & \cong
\mathcal{B}_{1}\left(  \mathcal{H}^{2}\right)
\end{align}
and hence%
\begin{align}
&  \dim\left(  \mathcal{K}_{31}\oplus\mathcal{K}_{30}\right)  =\dim\left(
\mathcal{B}_{1}\left(  \mathcal{H}^{1}\right)  \right)  =16\nonumber\\
&  \dim\left(  \mathcal{K}_{32}\oplus\mathcal{K}_{31}\oplus\mathcal{K}%
_{30}\right)  =\dim\left(  \mathcal{B}_{1}\left(  \mathcal{H}^{2}\right)
\right)  =36\\
&  \dim\left(  \mathcal{K}_{30}\right)  =1\nonumber\\
&  \dim\left(  \mathcal{K}_{32}\right)  =\dim\left(  \mathcal{B}_{1}\left(
\mathcal{H}^{2}\right)  \right)  -\dim\left(  \mathcal{B}_{1}\left(
\mathcal{H}^{1}\right)  \right)  =36-16=20\nonumber\\
&  \dim\left(  \mathcal{K}_{22od}\right)  =18\nonumber\\
&  \dim\left(  \mathcal{K}_{22d}\right)  =2\nonumber\\
&  \dim\left(  \mathcal{K}_{21od}\right)  =12\nonumber\\
&  \dim\left(  \mathcal{K}_{21d}\right)  =3
\end{align}


The space $\mathcal{K}_{1}=\ker L_{N}^{0}\ominus\ker L_{N}^{1}$ contains trace
class operators that have zero trace but have non zero one trace. This space
contains the spaces $\left\{  \mathfrak{B}_{N-1}\left(  \mathbf{\varphi
}\right)  \right\}  $ and $\mathfrak{B}_{N-1}\left(  \mathsf{z}\right)  $
i.e.$\mathcal{K}_{1}$ can be decomposed as the direct sum
\begin{equation}
\mathcal{K}_{1}=\mathfrak{B}_{N-1}\left(  \mathbf{\varphi}\right)
\oplus\mathfrak{B}_{N-1}\left(  \mathbf{\varphi}\right)  ^{c_{\mathcal{K}_{1}%
}}%
\end{equation}
so that
\begin{equation}
X_{1}=\sum_{\substack{1\leq j_{2}<\cdots<j_{N}\leq\infty\\1\leq j,k\leq\infty
}}X_{jkj_{2}\cdots j_{N}}\left\vert \varphi_{j}\varphi_{j_{2}}\cdots
\varphi_{j_{N}}\right\rangle \left\langle \varphi_{k}\varphi_{j_{2}}%
\cdots\varphi_{j_{N}}\right\vert ;\qquad L_{N}^{0}\left(  X\right)
=0;L_{N}^{1}\left(  X\right)  \neq0
\end{equation}
can be further develope as a sum of an \textit{off diagonal }and a
\textit{diagonal }component.\textit{ }
\begin{align}
X_{1} &  =\sum_{\substack{1\leq j_{2}<\cdots<j_{N}\leq\infty\\1\leq j\neq
k\leq\infty}}X_{jkj_{2}\cdots j_{N}}\left\vert \varphi_{j}\varphi_{j_{2}%
}\cdots\varphi_{j_{N}}\right\rangle \left\langle \varphi_{k}\varphi_{j_{2}%
}\cdots\varphi_{j_{N}}\right\vert \\
&  +\frac{1}{2}\sum_{\substack{1\leq j_{2}<\cdots<j_{N}\leq\infty\\1\leq
j\leq\infty}}X_{jj_{2}\cdots j_{N}}\left\{  \left\vert \varphi_{j}%
\varphi_{j_{2}}\cdots\varphi_{j_{N}}\right\rangle \left\langle \varphi
_{j}\varphi_{j_{2}}\cdots\varphi_{j_{N}}\right\vert -\left\vert \varphi
_{j+1}\varphi_{j_{2}}\cdots\varphi_{j_{N}}\right\rangle \left\langle
\varphi_{j+1}\varphi_{j_{2}}\cdots\varphi_{j_{N}}\right\vert \right\}  \\
&  \equiv X_{1K_{od}}+X_{1K_{d}}%
\end{align}
where
\begin{equation}
X_{1K_{od}}=\sum_{\substack{1\leq j_{2}<\cdots<j_{N}\leq\infty\\1\leq j\neq
k\leq\infty}}X_{jkj_{2}\cdots j_{N}}\left\vert \varphi_{j}\varphi_{j_{2}%
}\cdots\varphi_{j_{N}}\right\rangle \left\langle \varphi_{k}\varphi_{j_{2}%
}\cdots\varphi_{j_{N}}\right\vert \in\mathfrak{B}_{N-1}\left(  \mathbf{\varphi
}\right)
\end{equation}
and
\begin{equation}
X_{1K_{d}}=\frac{1}{2}\sum_{\substack{1\leq j_{2}<\cdots<j_{N}\leq
\infty\\1\leq j\leq\infty}}X_{jj_{2}\cdots j_{N}}\left\{  \left\vert
\varphi_{j}\varphi_{j_{2}}\cdots\varphi_{j_{N}}\right\rangle \left\langle
\varphi_{j}\varphi_{j_{2}}\cdots\varphi_{j_{N}}\right\vert -\left\vert
\varphi_{j+1}\varphi_{j_{2}}\cdots\varphi_{j_{N}}\right\rangle \left\langle
\varphi_{j+1}\varphi_{j_{2}}\cdots\varphi_{j_{N}}\right\vert \right\}
\in\mathfrak{B}_{N-1}\left(  \mathbf{\varphi}\right)  ^{c_{\mathcal{K}_{1}}}%
\end{equation}


The space $\mathcal{K}_{2}=\ker L_{N}^{1}\ominus\ker L_{N}^{2}$ contains trace
class operators that have zero one trace but have non zero two trace. This
space contains the spaces $\left\{  \mathfrak{B}_{N-2}\left(  \mathbf{\varphi
}\right)  \right\}  $ and $\mathfrak{B}_{N-2}\left(  \mathsf{z}\right)  $
i.e.$\mathcal{K}_{2}$ can be decomposed as the direct sum
\begin{equation}
\mathcal{K}_{2}=\mathfrak{B}_{N-2}\left(  \mathbf{\varphi}\right)
\oplus\mathfrak{B}_{N-2}\left(  \mathbf{\varphi}\right)  ^{c_{\mathcal{K}_{2}%
}}%
\end{equation}
so that
\begin{equation}
X_{2}=\sum_{\substack{1\leq j_{3}<\cdots<j_{N}\leq\infty\\1\leq j_{1}%
<,j_{2},k_{1}<k_{2}\leq\infty}}X_{j_{1}j_{2}k_{1}k_{2}j_{3}\cdots j_{N}%
}\left\vert \varphi_{j_{1}}\varphi_{j_{2}}\varphi_{j_{3}}\cdots\varphi_{j_{N}%
}\right\rangle \left\langle \varphi_{k_{1}}\varphi_{k_{2}}\varphi_{j_{3}%
}\cdots\varphi_{j_{N}}\right\vert ;\qquad L_{N}^{1}\left(  X\right)
=0;L_{N}^{2}\left(  X\right)  \neq0
\end{equation}
can be further developed as a sum of an \textit{off diagonal }and a
\textit{diagonal }component.\textit{ }
\begin{align}
&  X_{2}=\sum_{\substack{1\leq j_{3}<\cdots<j_{N}\leq\infty\\1\leq
j_{1}<,j_{2}\neq k_{1}<k_{2}\leq\infty}}X_{j_{1}j_{2}k_{1}k_{2}j_{3}\cdots
j_{N}}\left\vert \varphi_{j_{1}}\varphi_{j_{2}}\varphi_{j_{3}}\cdots
\varphi_{j_{N}}\right\rangle \left\langle \varphi_{k_{1}}\varphi_{k_{2}%
}\varphi_{j_{3}}\cdots\varphi_{j_{N}}\right\vert +\\
&  \frac{1}{2}\sum_{\substack{1\leq j_{3}<\cdots<j_{N}\leq\infty\\1\leq
j_{1}\neq k_{1}\leq\infty}}X_{j_{1}k_{1}j_{2}j_{3}\cdots j_{N}}\left\{
\left\vert \varphi_{j_{1}}\varphi_{j_{2}}\varphi_{j_{3}}\cdots\varphi_{j_{N}%
}\right\rangle \left\langle \varphi_{k_{1}}\varphi_{j_{2}}\varphi_{j_{3}%
}\cdots\varphi_{j_{N}}\right\vert -\left\vert \varphi_{j_{1}}\varphi_{j_{2}%
+1}\varphi_{j_{3}}\cdots\varphi_{j_{N}}\right\rangle \left\langle
\varphi_{k_{1}}\varphi_{j_{2}+1}\varphi_{j_{3}}\cdots\varphi_{j_{N}%
}\right\vert \right\} \\
&  \equiv X_{2K_{od}}+X_{2K_{d}}%
\end{align}
where
\begin{equation}
X_{1K_{od}}=\sum_{\substack{1\leq j_{2}<\cdots<j_{N}\leq\infty\\1\leq j\neq
k\leq\infty}}X_{jkj_{2}\cdots j_{N}}\left\vert \varphi_{j}\varphi_{j_{2}%
}\cdots\varphi_{j_{N}}\right\rangle \left\langle \varphi_{k}\varphi_{j_{2}%
}\cdots\varphi_{j_{N}}\right\vert \in\mathfrak{B}_{N-1}\left(  \mathbf{\varphi
}\right)
\end{equation}
and
\begin{equation}
X_{1K_{d}}=\frac{1}{2}\sum_{\substack{1\leq j_{2}<\cdots<j_{N}\leq
\infty\\1\leq j\leq\infty}}X_{jj_{2}\cdots j_{N}}\left\{  \left\vert
\varphi_{j}\varphi_{j_{2}}\cdots\varphi_{j_{N}}\right\rangle \left\langle
\varphi_{j}\varphi_{j_{2}}\cdots\varphi_{j_{N}}\right\vert -\left\vert
\varphi_{j+1}\varphi_{j_{2}}\cdots\varphi_{j_{N}}\right\rangle \left\langle
\varphi_{j+1}\varphi_{j_{2}}\cdots\varphi_{j_{N}}\right\vert \right\}
\in\mathfrak{B}_{N-1}\left(  \mathbf{\varphi}\right)  ^{c}%
\end{equation}
The

In particular one can use the basis $\mathfrak{v}_{2}$ of $\mathfrak{V}_{2}$
to construct a new mixed generalized-proper basis of $\mathcal{B}_{1}\left(
\mathcal{H}^{N}\right)  $ that is adapted to the decomposition eq. $\left(
\ref{orthog}\right)  $. There are three types of elements in this adapted basis:

\begin{enumerate}
\item The generalized operators $\mathfrak{v}_{2}=\left\{  \left\vert
\mathsf{z}^{N}\right\rangle \left\langle \mathsf{z}^{\prime N}\right\vert
\left\vert \mathsf{z}^{N}\overset{2}{\nsim}\mathsf{z}^{\prime N}\right.
\right\}  $. These span the space $\mathfrak{V}_{2}$ but which we will denote
as $\mathcal{K}_{OD}$ in this context. This basis could be changed into a
proper basis by defining trace class "wavepackets"
\[
l_{\mu}=%
%TCIMACRO{\dint }%
%BeginExpansion
{\displaystyle\int}
%EndExpansion
l_{\mu}\left(  \mathsf{z}^{N};\mathsf{z}^{\prime N}\right)  \left\vert
\mathsf{z}^{N}\right\rangle \left\langle \mathsf{z}^{\prime N}\right\vert
d\mathsf{z}^{N}d\mathsf{z}^{\prime N};\qquad\mathsf{z}^{N}\overset{2}{\nsim
}\mathsf{z}^{\prime N}%
\]
but for the purposes of this article we find it more convenient to stay with
the generalized basis.

\item The proper trace class operators that are elements of $\ker L_{N}^{1}$
\begin{align*}
&  \left\{  k_{\kappa}=%
%TCIMACRO{\dint }%
%BeginExpansion
{\displaystyle\int}
%EndExpansion
k_{\kappa}\left(  \mathsf{z}_{1}\mathsf{z}^{N-1};\mathsf{z}_{1}^{\prime
}\mathsf{z}^{N-1}\right)  \left\vert \mathsf{z}_{1}\mathsf{z}^{N-1}%
\right\rangle \left\langle \mathsf{z}_{1}^{\prime}\mathsf{z}^{N-1}\right\vert
d\mathsf{z}_{1}d\mathsf{z}_{1}^{\prime}d\mathsf{z}^{N-1}\right. \\
&  \qquad\qquad\left.  \left\vert
%TCIMACRO{\dint }%
%BeginExpansion
{\displaystyle\int}
%EndExpansion
k_{\kappa}\left(  \mathsf{z}_{1}\mathsf{z}^{N-1};\mathsf{z}_{1}^{\prime
}\mathsf{z}^{N-1}\right)  d\mathsf{z}^{N-1}\right.  \overset{a.e}{=}%
0\quad\forall\mathsf{z}_{1}\mathsf{z}_{1}^{\prime};1\leq\kappa\leq
\infty\overset{}{}\right\}
\end{align*}


which span a space which we will denote $\mathcal{K}_{D}$.

\item The proper trace class operators that are elements of $\left[  \ker
L_{N}^{1}\right]  ^{c}$%
\begin{align*}
&  \left\{  \mathfrak{x}_{\lambda}=%
%TCIMACRO{\dint }%
%BeginExpansion
{\displaystyle\int}
%EndExpansion
\mathfrak{x}_{\lambda}\left(  \mathsf{z}_{1}\mathsf{z}^{N-1};\mathsf{z}%
_{1}^{\prime}\mathsf{z}^{N-1}\right)  \left\vert \mathsf{z}_{1}\mathsf{z}%
^{N-1}\right\rangle \left\langle \mathsf{z}_{1}^{\prime}\mathsf{z}%
^{N-1}\right\vert d\mathsf{z}_{1}d\mathsf{z}_{1}^{\prime}d\mathsf{z}%
^{N-1}|\right. \\
&  \qquad\qquad\qquad\left.  \int\mathfrak{x}_{\lambda}\left(  \mathsf{z}%
_{1}\mathsf{z}^{N-1};\mathsf{z}_{1}^{\prime}\mathsf{z}^{N-1}\right)
d\mathsf{z}^{N-1}\neq0\quad\forall\mathsf{z}_{1}\mathsf{z}_{1}^{\prime}%
\in\Delta;1\leq\lambda\leq\infty\overset{}{}\right\}
\end{align*}


where $\Delta$ is a set of non zero measure. These span the space $\left[
\ker L_{N}^{1}\right]  ^{c}$.
\end{enumerate}

The basis elements of type 1 and type 2 together clearly span $\ker L_{N}^{1}
$, which can be expressed as the direct sum%

\begin{equation}
\ker L_{N}^{1}=\mathcal{K}_{OD}\oplus\mathcal{K}_{D}.
\end{equation}
$\mathcal{K}_{OD}$ can be further decomposed as direct a integral
\cite{dirint} of subspaces, $\mathcal{K}_{\left\vert \mathsf{z}_{1}%
\mathsf{z}^{N-1}\right\rangle \left\langle \mathsf{z}_{1}^{\prime}%
\mathsf{z}^{\prime N}\right\vert }$, indexed by the generalized basis elements
$\left\{  \left\vert \mathsf{z}^{N}\right\rangle \left\langle \mathsf{z}%
^{\prime N}\right\vert \right\}  $, where $\mathsf{z}^{N}\overset{2}{\nsim
}\mathsf{z}^{\prime N}$ i.e.
\begin{equation}
\mathcal{K}_{OD}\equiv\mathfrak{V}_{2}=\int_{\oplus}\mathcal{K}_{\left\vert
\mathsf{z}^{N}\right\rangle \left\langle \mathsf{z}^{\prime N}\right\vert
}d\mathsf{z}^{N}d\mathsf{z}^{\prime N};\quad\mathsf{z}^{N}\overset{2}{\nsim
}\mathsf{z}^{\prime N}.
\end{equation}


One can define a basis of the dual space $\mathcal{B}\left(  \mathcal{H}%
^{N}\right)  $ that is biorthogonal to the basis of $\mathcal{B}_{1}\left(
\mathcal{H}^{N}\right)  $ just described by%
\begin{equation}
\left\{  \left\{  \left\vert \mathsf{z}^{N}\right\rangle \left\langle
\mathsf{z}^{\prime N}\right\vert \left\vert \mathsf{z}^{N}\overset{2}{\nsim
}\mathsf{z}^{\prime N}\right.  \right\}  ,\left\{  k_{\kappa}^{d}\right\}
,\left\{  \mathfrak{x}_{\lambda}^{d}\right\}  \right\}
\end{equation}
where
\begin{align}
\left(  k_{\kappa}^{d}|k_{\lambda}\right)   &  =\delta_{\kappa\lambda
}\nonumber\\
\left(  \mathfrak{x}_{\kappa}^{d}|\mathfrak{x}_{\lambda}\right)   &
=\delta_{\kappa\lambda}%
\end{align}
and
\begin{equation}
\left(  k_{\kappa}^{d}|b\right)  =\left(  \mathfrak{x}_{\kappa}^{d}|b\right)
=0
\end{equation}
for all other basis elements $\left\{  \left\vert b\right)  \right\}  $ of
$\mathcal{B}_{1}\left(  \mathcal{H}^{N}\right)  .$ The dual basis elements of
$\mathcal{B}\left(  \mathcal{H}^{N}\right)  $ are explicitly expressed as:

\begin{enumerate}
\item The same generalized operators $\left\{  \left\vert \mathsf{z}%
^{N}\right\rangle \left\langle \mathsf{z}^{\prime N}\right\vert \left\vert
\mathsf{z}^{N}\overset{2}{\nsim}\mathsf{z}^{\prime N}\right.  \right\}  $.
These span a space which we will denote as $\mathcal{K}_{OD}^{d}$. This basis
could be changed into a proper basis by defining bounded operator
"wavepackets"
\begin{equation}
l_{\mu}^{d}=%
%TCIMACRO{\dint }%
%BeginExpansion
{\displaystyle\int}
%EndExpansion
l_{\mu}^{d}\left(  \mathsf{z}^{N};\mathsf{z}^{\prime N}\right)  \left\vert
\mathsf{z}^{N}\right\rangle \left\langle \mathsf{z}^{\prime N}\right\vert
d\mathsf{z}^{N}d\mathsf{z}^{\prime N};\qquad\mathsf{z}^{N}\overset{2}{\nsim
}\mathsf{z}^{\prime N}%
\end{equation}
but again for the purposes of this article we find it more convenient to stay
with the generalized basis.

\item The proper bounded operators%
\begin{align}
&  \left\{  k_{\kappa}^{d}=%
%TCIMACRO{\dint }%
%BeginExpansion
{\displaystyle\int}
%EndExpansion
k_{\kappa}^{d}\left(  \mathsf{z}_{1}\mathsf{z}^{N-1};\mathsf{z}_{1}^{\prime
}\mathsf{z}^{N-1}\right)  \left\vert \mathsf{z}_{1}\mathsf{z}^{N-1}%
\right\rangle \left\langle \mathsf{z}_{1}^{\prime}\mathsf{z}^{N-1}\right\vert
d\mathsf{z}_{1}d\mathsf{z}_{1}^{\prime}d\mathsf{z}^{N-1}\right. \nonumber\\
&  \qquad\qquad\qquad\left.  \left\vert
%TCIMACRO{\dint }%
%BeginExpansion
{\displaystyle\int}
%EndExpansion
k_{\kappa}^{d}\left(  \mathsf{z}_{1}\mathsf{z}^{N-1};\mathsf{z}_{1}^{\prime
}\mathsf{z}^{N-1}\right)  d\mathsf{z}^{N-1}\right.  \overset{a.e}{=}%
0\quad\forall\mathsf{z}_{1}\mathsf{z}_{1}^{\prime};1\leq\kappa\leq
\infty\overset{}{}\right\}
\end{align}


which span a space which we will denote $\mathcal{K}_{D}^{d}$.

\item The proper bounded operators%
\begin{align}
&  \left\{  \mathfrak{x}_{\lambda}^{d}=%
%TCIMACRO{\dint }%
%BeginExpansion
{\displaystyle\int}
%EndExpansion
\mathfrak{x}_{\lambda}^{d}\left(  \mathsf{z}_{1}\mathsf{z}^{N-1}%
;\mathsf{z}_{1}^{\prime}\mathsf{z}^{N-1}\right)  \left\vert \mathsf{z}%
_{1}\mathsf{z}^{N-1}\right\rangle \left\langle \mathsf{z}_{1}^{\prime
}\mathsf{z}^{N-1}\right\vert d\mathsf{z}_{1}d\mathsf{z}_{1}^{\prime
}d\mathsf{z}^{N-1}|\right. \nonumber\\
&  \qquad\qquad\qquad\left.  \int\mathfrak{x}_{\lambda}^{d}\left(
\mathsf{z}_{1}\mathsf{z}^{N-1};\mathsf{z}_{1}^{\prime}\mathsf{z}^{N-1}\right)
d\mathsf{z}^{N-1}\neq0\quad\forall\mathsf{z}_{1}\mathsf{z}_{1}^{\prime}%
\in\Delta;1\leq\lambda\leq\infty\overset{}{}\right\}
\end{align}


where $\Delta$ is a set of non zero measure. These span the space $\left[
\ker L_{N}^{1}\right]  ^{cd}$. It \emph{should} be noted that in general%
\begin{equation}
k_{\kappa}^{d}\left(  \mathsf{z}_{1}\mathsf{z}^{N-1};\mathsf{z}_{1}^{\prime
}\mathsf{z}^{N-1}\right)  \neq k_{\kappa}\left(  \mathsf{z}_{1}\mathsf{z}%
^{N-1};\mathsf{z}_{1}^{\prime}\mathsf{z}^{N-1}\right)  ^{\ast}%
\end{equation}
and
\begin{equation}
\mathfrak{x}_{\lambda}^{d}\left(  \mathsf{z}_{1}\mathsf{z}^{N-1}%
;\mathsf{z}_{1}^{\prime}\mathsf{z}^{N-1}\right)  \neq\mathfrak{x}_{\lambda
}\left(  \mathsf{z}_{1}\mathsf{z}^{N-1};\mathsf{z}_{1}^{\prime}\mathsf{z}%
^{N-1}\right)  ^{\ast}.
\end{equation}

\end{enumerate}

The biorthogonality of the preceding bases can be displayed symbolically as%
\begin{equation}
\left(
\begin{array}
[c]{c}%
\left(  \underline{\left\vert \mathsf{z}^{N}\right\rangle _{2}\left\langle
\mathsf{z}^{\prime N}\right\vert }\right\vert \\
\left(  \underline{k^{d}}\right\vert \\
\left(  \underline{\mathfrak{x}^{d}}\right\vert
\end{array}
\right)  \left(  \left\vert \underline{\left\vert \mathsf{z}^{N}\right\rangle
_{2}\left\langle \mathsf{z}^{\prime N}\right\vert }\right)  ,\underline
{\left\vert k\right)  },\underline{\left\vert \mathfrak{x}^{d}\right)
}\right)  =\left(
\begin{array}
[c]{ccc}%
\mathbb{I}_{\mathcal{K}_{OD}} & \mathbb{O} & \mathbb{O}\\
\mathbb{O} & \mathbb{I}_{\mathcal{K}_{D}} & \mathbb{O}\\
\mathbb{O} & \mathbb{O} & \mathbb{I}_{\left[  \ker L_{N}^{1}\right]  ^{c}}%
\end{array}
\right)
\end{equation}
This basis also provides us with a resolution of the identity symbolically
expressed as%
\begin{equation}
\hat{I}_{\mathcal{B}_{1}\left(  \mathcal{H}^{N}\right)  }=\left(
\underline{\left\vert \left\vert \mathsf{z}^{N}\right\rangle _{2}\left\langle
\mathsf{z}^{\prime N}\right\vert \right)  },\underline{\left\vert k\right)
},\underline{\left\vert \mathfrak{x}\right)  }\right)  \left(
\begin{array}
[c]{c}%
\underline{\left(  \left\vert \mathsf{z}^{N}\right\rangle _{2}\left\langle
\mathsf{z}^{\prime N}\right\vert \right\vert }\\
\underline{\left(  k^{d}\right\vert }\\
\underline{\left(  \mathfrak{x}^{d}\right\vert }%
\end{array}
\right)  \label{rescan}%
\end{equation}
where $\left\vert \left\vert \mathsf{z}^{N}\right\rangle _{2}\left\langle
\mathsf{z}^{\prime N}\right\vert \right)  $ denotes the basis $\left\{
\left\vert \mathsf{z}^{N}\right\rangle \left\langle \mathsf{z}^{\prime
N}\right\vert \left\vert \mathsf{z}^{N}\overset{2}{\nsim}\mathsf{z}^{\prime
N}\right.  \right\}  $.

Any element $X\in\mathcal{B}_{1}\left(  \mathcal{H}^{N}\right)  $ can thus be
expanded as
\begin{equation}
X=X_{\mathcal{K}_{OD}}+X_{\mathcal{K}_{D}}+X_{\left[  \ker L_{N}^{1}\right]
^{c}} \label{kerdec}%
\end{equation}
where
\begin{align}
X_{\mathcal{K}_{OD}}  &  =\int\left(  \left\vert \mathsf{z}^{N}\right\rangle
\left\langle \mathsf{z}^{\prime N}\right\vert |X\right)  \left\vert
\mathsf{z}^{N}\right\rangle \left\langle \mathsf{z}^{\prime N}\right\vert
d\mathsf{z}^{N}d\mathsf{z}^{\prime N}\nonumber\\
&  \int Tr\left\{  \left\vert \mathsf{z}^{N}\right\rangle \left\langle
\mathsf{z}^{\prime N}\right\vert ^{\dagger}X\right\}  \left\vert
\mathsf{z}^{N}\right\rangle \left\langle \mathsf{z}^{\prime N}\right\vert
d\mathsf{z}^{N}d\mathsf{z}^{\prime N}\nonumber\\
&  =\int\left\langle \mathsf{z}^{N}|X\mathsf{z}^{\prime N}\right\rangle
\left\vert \mathsf{z}^{N}\right\rangle \left\langle \mathsf{z}^{\prime
N}\right\vert d\mathsf{z}^{N}d\mathsf{z}^{\prime N}\nonumber\\
&  =\int X\left(  \mathsf{z}^{N},\mathsf{z}^{\prime N}\right)  \left\vert
\mathsf{z}^{N}\right\rangle \left\langle \mathsf{z}^{\prime N}\right\vert
d\mathsf{z}^{N}d\mathsf{z}^{\prime N};\quad\mathsf{z}^{N}\overset{2}{\nsim
}\mathsf{z}^{\prime N}%
\end{align}%
\begin{equation}
\qquad X_{\mathcal{K}_{D}}=%
%TCIMACRO{\dsum \limits_{\kappa}}%
%BeginExpansion
{\displaystyle\sum\limits_{\kappa}}
%EndExpansion
\left(  k_{\kappa}^{d}|X\right)  k_{\kappa}=%
%TCIMACRO{\dsum \limits_{\kappa}}%
%BeginExpansion
{\displaystyle\sum\limits_{\kappa}}
%EndExpansion
Tr\left\{  k_{\kappa}^{d\dagger}X\right\}  k_{\kappa}\qquad\qquad\qquad
\end{equation}
and
\begin{equation}
X_{\left[  \ker L_{N}^{1}\right]  ^{c}}=%
%TCIMACRO{\dsum \limits_{\kappa}}%
%BeginExpansion
{\displaystyle\sum\limits_{\kappa}}
%EndExpansion
\left(  \mathfrak{x}_{\lambda}^{d}|X\right)  \mathfrak{x}_{\lambda}=%
%TCIMACRO{\dsum \limits_{\kappa}}%
%BeginExpansion
{\displaystyle\sum\limits_{\kappa}}
%EndExpansion
Tr\left\{  \mathfrak{x}_{\lambda}^{d\dagger}X\right\}  \mathfrak{x}_{\lambda
}\qquad\qquad\qquad
\end{equation}
When analyzing the properties of an operator $X\in\mathcal{B}_{1}\left(
\mathcal{H}^{N}\right)  $ it behooves us to both consider its expansion wrt
the subspace decomposition $\mathcal{B}_{1}\left(  \mathcal{H}^{N}\right)  $
induced by $\ker L_{N}^{1}$ as in eq. $\left(  \ref{kerdec}\right)  $ and its
expansion wrt to subspace decomposition of $\mathcal{B}_{1}\left(
\mathcal{H}^{N}\right)  $ given by $\left\{  \mathfrak{B}_{m}|0\leq m\leq
N\right\}  $ i.e.
\begin{equation}
X=\sum_{0\leq m\leq N}X_{\mathfrak{B}_{m}}%
\end{equation}
where
\begin{equation}
X_{\mathfrak{B}_{m}}=%
%TCIMACRO{\dint \limits_{\mathsf{z}^{\prime N}\,\overset{m}{=}\,\mathsf{z}^{N}%
%}}%
%BeginExpansion
{\displaystyle\int\limits_{\mathsf{z}^{\prime N}\,\overset{m}{=}%
\,\mathsf{z}^{N}}}
%EndExpansion
X_{m}\left(  \mathsf{z}^{\prime N},\mathsf{z}^{N}\right)  \left\vert
\mathsf{z}^{\prime N}\right\rangle \left\langle \mathsf{z}^{N}\right\vert
d\mathsf{z}^{N}d\mathsf{z}^{\prime N}%
\end{equation}
Explicitly using the resolution of the identity for the basis $\left\{
\underline{\left\vert \left\vert \mathsf{z}^{N}\right\rangle _{2}\left\langle
\mathsf{z}^{\prime N}\right\vert \right)  },\underline{\left\vert k\right)
},\underline{\left\vert \mathfrak{x}\right)  }\right\}  $ we obtain%
\begin{align}
X  &  =\left(  \underline{\left\vert \left\vert \mathsf{z}^{N}\right\rangle
_{2}\left\langle \mathsf{z}^{\prime N}\right\vert \right)  },\underline
{\left\vert k\right)  },\underline{\left\vert \mathfrak{x}\right)  }\right)
\left(
\begin{array}
[c]{c}%
\underline{\left(  \left\vert \mathsf{z}^{N}\right\rangle _{2}\left\langle
\mathsf{z}^{\prime N}\right\vert \right\vert }\\
\underline{\left(  k^{d}\right\vert }\\
\underline{\left(  \mathfrak{x}^{d}\right\vert }%
\end{array}
\right)  \left\vert X\right) \nonumber\\
&  =\left(  \underline{\left\vert \left\vert \mathsf{z}^{N}\right\rangle
_{2}\left\langle \mathsf{z}^{\prime N}\right\vert \right)  },\underline
{\left\vert k\right)  },\underline{\left\vert \mathfrak{x}\right)  }\right)
\mathbb{X}_{can}%
\end{align}
and for the basis $\left\{  \mathfrak{b}_{0},\cdots,\mathfrak{b}_{N}\right\}
$%
\begin{equation}
X=\left(  \mathfrak{b}_{0},\cdots,\mathfrak{b}_{N}\right)  \left(
\begin{array}
[c]{c}%
\mathfrak{b}_{0}\\
\vdots\\
\mathfrak{b}_{N}%
\end{array}
\right)  \left\vert X\right)  =\left(  \mathfrak{b}_{0},\cdots,\mathfrak{b}%
_{N}\right)  \mathbb{X}_{\mathcal{B}_{1}\left(  \mathcal{H}^{N}\right)  }%
\end{equation}
leading to the transformation from the representation $\mathbb{X}_{can}$ to
$\mathbb{X}_{\mathcal{B}_{1}\left(  \mathcal{H}^{N}\right)  }$ given as
\begin{align}
\mathbb{X}_{can}  &  =\left(
\begin{array}
[c]{c}%
\underline{\left(  \left\vert \mathsf{z}^{N}\right\rangle _{2}\left\langle
\mathsf{z}^{\prime N}\right\vert \right\vert }\\
\underline{\left(  k^{d}\right\vert }\\
\underline{\left(  \mathfrak{x}^{d}\right\vert }%
\end{array}
\right)  \left(  \mathfrak{b}_{0},\cdots,\mathfrak{b}_{N}\right)
\mathbb{X}_{\mathcal{B}_{1}\left(  \mathcal{H}^{N}\right)  }\nonumber\\
&  =\left(
\begin{array}
[c]{ccc}%
\mathbb{I} & \mathbb{O} & \mathbb{O}\\
\mathbb{O} & \left(  \left.  \underline{\left(  k^{d}\right\vert }\right\vert
\mathfrak{b}_{N-1}\right)  & \left(  \left.  \underline{\left(  k^{d}%
\right\vert }\right\vert \mathfrak{b}_{N}\right) \\
\mathbb{O} & \left(  \left.  \underline{\left(  \mathfrak{x}^{d}\right\vert
}\right\vert \mathfrak{b}_{N-1}\right)  & \left(  \left.  \underline{\left(
\mathfrak{x}^{d}\right\vert }\right\vert \mathfrak{b}_{N}\right)
\end{array}
\right)  \mathbb{X}_{\mathcal{B}_{1}\left(  \mathcal{H}^{N}\right)  }%
\end{align}
and the transformation from the representation $\mathbb{X}_{\mathcal{B}%
_{1}\left(  \mathcal{H}^{N}\right)  }$ to $\mathbb{X}_{can}$ as
\begin{align}
\mathbb{X}_{\mathcal{B}_{1}\left(  \mathcal{H}^{N}\right)  }  &  =\left(
\begin{array}
[c]{c}%
\mathfrak{b}_{0}\\
\vdots\\
\mathfrak{b}_{N}%
\end{array}
\right)  \left(  \underline{\left\vert \left\vert \mathsf{z}^{N}\right\rangle
_{2}\left\langle \mathsf{z}^{\prime N}\right\vert \right)  },\underline
{\left\vert k\right)  },\underline{\left\vert \mathfrak{x}\right)  }\right)
\mathbb{X}_{can}\nonumber\\
&  =\left(
\begin{array}
[c]{ccc}%
\mathbb{I} & \mathbb{O} & \mathbb{O}\\
\mathbb{O} & \left(  \mathfrak{b}_{N-1}\left\vert \underline{\left\vert
k\right)  }\right.  \right)  & \left(  \mathfrak{b}_{N-1}\left\vert
\underline{\left\vert \mathfrak{x}\right)  }\right.  \right) \\
\mathbb{O} & \left(  \mathfrak{b}_{N}\left\vert \underline{\left\vert
k\right)  }\right.  \right)  & \left(  \mathfrak{b}_{N}\left\vert
\underline{\left\vert \mathfrak{x}\right)  }\right.  \right)
\end{array}
\right)  \mathbb{X}_{can}%
\end{align}
Explicitly the expansion coefficients that change are given by
\begin{subequations}
\begin{align}
X_{k_{j}}  &  =%
%TCIMACRO{\dint \limits_{\mathsf{z}_{1}^{N}\,\overset{_{1}}{=}\,\mathsf{z}%
%_{2}^{N};\mathsf{z}_{1}^{N}\,\overset{}{=}\,\mathsf{z}_{2}^{N}}}%
%BeginExpansion
{\displaystyle\int\limits_{\mathsf{z}_{1}^{N}\,\overset{_{1}}{=}%
\,\mathsf{z}_{2}^{N};\mathsf{z}_{1}^{N}\,\overset{}{=}\,\mathsf{z}_{2}^{N}}}
%EndExpansion
k_{j}^{d}\left(  \mathsf{z}^{N},\mathsf{z}^{\prime N}\right)  ^{\ast}X\left(
\mathsf{z}^{N},\mathsf{z}^{\prime N}\right)  d\mathsf{z}_{1}^{N}%
d\mathsf{z}_{2}^{N}\\
X_{\mathfrak{x}_{j}}  &  =%
%TCIMACRO{\dint \limits_{\mathsf{z}_{1}^{N}\,\overset{_{1}}{=}\,\mathsf{z}%
%_{2}^{N};\mathsf{z}_{1}^{N}\,\overset{}{=}\,\mathsf{z}_{2}^{N}}}%
%BeginExpansion
{\displaystyle\int\limits_{\mathsf{z}_{1}^{N}\,\overset{_{1}}{=}%
\,\mathsf{z}_{2}^{N};\mathsf{z}_{1}^{N}\,\overset{}{=}\,\mathsf{z}_{2}^{N}}}
%EndExpansion
\mathfrak{x}_{j}^{d}\left(  \mathsf{z}^{N},\mathsf{z}^{\prime N}\right)
^{\ast}X\left(  \mathsf{z}^{N},\mathsf{z}^{\prime N}\right)  d\mathsf{z}%
_{1}^{N}d\mathsf{z}_{2}^{N}%
\end{align}
and
\end{subequations}
\begin{equation}
X\left(  \mathsf{z}^{N},\mathsf{z}^{\prime N}\right)  =%
%TCIMACRO{\dsum \limits_{1\leq j\leq\infty}}%
%BeginExpansion
{\displaystyle\sum\limits_{1\leq j\leq\infty}}
%EndExpansion
\left\{  X_{k_{j}}k_{j}\left(  \mathsf{z}^{N},\mathsf{z}^{\prime N}\right)
+X_{\mathfrak{x}_{j}}\mathfrak{x}_{j}\left(  \mathsf{z}^{N},\mathsf{z}^{\prime
N}\right)  \right\}  ;\qquad\mathsf{z}_{1}^{N}\,\overset{_{1}}{=}%
\,\mathsf{z}_{2}^{N};\mathsf{z}_{1}^{N}\,\overset{}{=}\,\mathsf{z}_{2}^{N}%
\end{equation}


\subsubsection{Non Canonical Basis of Complimentary Subspaces of $\ker
L_{N}^{1}.$}

We can utilize the maps $\tilde{\Gamma}$ introduced in eq. $\left(
\ref{Csubspace}\right)  $ to define a map $\Upsilon:\mathcal{B}_{1}\left(
\mathcal{H}^{N}\right)  \rightarrow\mathcal{B}_{1}\left(  \mathcal{H}%
^{N}\right)  $ to produce bases adapted to the decompositions eq. $\left(
\ref{nonorthog}\right)  $ as
\begin{align}
&  \left\{  \left\{  \left\vert \mathsf{z}^{N}\right\rangle \left\langle
\mathsf{z}^{\prime N}\right\vert \,\forall\mathsf{z}^{N}\overset{2}{\nsim
}\mathsf{z}^{\prime N}\right\}  ,\left\{  k_{\kappa}\,;1\leq\lambda\leq
\infty\right\}  ,\left\{  \tilde{\Gamma}\left(  \mathfrak{x}_{\lambda}\right)
\,;1\leq\lambda\leq\infty\right\}  \right\} \nonumber\\
&  \equiv\left\{  \left\{  \Upsilon\left(  \left\vert \mathsf{z}%
^{N}\right\rangle \left\langle \mathsf{z}^{\prime N}\right\vert \right)
\,\forall\mathsf{z}^{N}\overset{2}{\nsim}\mathsf{z}^{\prime N}\right\}
,\left\{  \Upsilon\left(  k_{\kappa}\right)  \,;1\leq\lambda\leq
\infty\right\}  ,\left\{  \Upsilon\left(  \mathfrak{x}_{\lambda}\right)
\,;1\leq\lambda\leq\infty\right\}  \right\}  \label{noncanbasis}%
\end{align}
where one should note that it is only the bases elements of $\left[  \ker
L_{N}^{1}\right]  ^{c}$ that have been changed. The\emph{\ }dual basis
elements $\left\{  \left(  k_{\kappa}^{d}\right\vert \right\}  $ are now no
longer biorthogonal to $\left\{  \left\vert \tilde{\Gamma}\left(
\mathfrak{x}_{\lambda}\right)  \right)  \right\}  $ as%
\begin{equation}
\tilde{\Gamma}\left(  \mathfrak{x}_{\lambda}\right)  =\Gamma\left(
\mathfrak{x}_{\lambda}\right)  +\mathfrak{x}_{\lambda}.
\end{equation}
and
\begin{equation}
\Gamma\left(  \mathfrak{x}_{\lambda}\right)  =\sum_{\kappa}\left\vert
k_{\kappa}\right)  \left[  \Gamma_{\mathcal{K}_{D}}\right]  _{\kappa\lambda}+%
%TCIMACRO{\dint _{\mathsf{z}^{N}\,\overset{2}{\neq\,}\mathsf{z}^{\prime N}}}%
%BeginExpansion
{\displaystyle\int_{\mathsf{z}^{N}\,\overset{2}{\neq\,}\mathsf{z}^{\prime N}}}
%EndExpansion
\left[  \Gamma_{\mathcal{K}_{OD}}\right]  _{\lambda}\left(  \mathsf{z}%
^{N};\mathsf{z}^{\prime N}\right)  \left\vert \mathsf{z}^{N}\right\rangle
\left\langle \mathsf{z}^{\prime N}\right\vert d\mathsf{z}^{N}d\mathsf{z}%
^{\prime N}%
\end{equation}
However one can construct a basis that is biorthogonal to eq. $\left(
\ref{noncanbasis}\right)  $ by considering
\begin{align}
&  \left(  \Upsilon_{\Gamma}^{d}\left(  \underline{\left(  \left\vert
\mathsf{z}^{N}\right\rangle _{2}\left\langle \mathsf{z}^{\prime N}\right\vert
\right\vert },\underline{\left(  k^{d}\right\vert },\left(  \underline
{\mathfrak{x}^{d}}\right\vert \right)  \right)  \left(  \Upsilon_{\Gamma
}\left(  \underline{\left\vert \left\vert \mathsf{z}^{N}\right\rangle
_{2}\left\langle \mathsf{z}^{\prime N}\right\vert \right)  },\underline
{\left\vert k\right)  },\underline{\left\vert \mathfrak{x}\right)  }\right)
\right)  ^{t}=\nonumber\\
&  \left(  \underline{\left(  \left\vert \mathsf{z}^{N}\right\rangle
_{2}\left\langle \mathsf{z}^{\prime N}\right\vert \right\vert },\underline
{\left(  k^{d}\right\vert },\left(  \underline{\mathfrak{x}^{d}}\right\vert
\right)  \left(
\begin{array}
[c]{ccc}%
\mathbb{I} & \mathbb{O} & -\mathbf{\Gamma}_{\mathcal{K}_{OD}}\\
\mathbb{O} & \mathbb{I} & -\mathbf{\Gamma}_{\mathcal{K}_{D}}\\
\mathbb{O} & \mathbb{O} & \mathbb{I}%
\end{array}
\right)  \left(
\begin{array}
[c]{ccc}%
\mathbb{I} & \mathbb{O} & \mathbf{\Gamma}_{\mathcal{K}_{OD}}\\
\mathbb{O} & \mathbb{I} & \mathbf{\Gamma}_{\mathcal{K}_{D}}\\
\mathbb{O} & \mathbb{O} & \mathbb{I}%
\end{array}
\right)  \left(
\begin{array}
[c]{c}%
\underline{\left\vert \left\vert \mathsf{z}^{N}\right\rangle _{2}\left\langle
\mathsf{z}^{\prime N}\right\vert \right)  }\\
\underline{\left\vert k\right)  }\\
\underline{\left\vert \mathfrak{x}\right)  }%
\end{array}
\right) \nonumber\\
&  \left.  =\right.  \left(  \underline{\left(  \left\vert \mathsf{z}%
^{N}\right\rangle _{2}\left\langle \mathsf{z}^{\prime N}\right\vert
\right\vert },\underline{\left(  k^{d}\right\vert },\left(  \underline
{\mathfrak{x}^{d}}\right\vert \right)  \left(
\begin{array}
[c]{c}%
\underline{\left\vert \left\vert \mathsf{z}^{N}\right\rangle _{2}\left\langle
\mathsf{z}^{\prime N}\right\vert \right)  }\\
\underline{\left\vert k\right)  }\\
\underline{\left\vert \mathfrak{x}\right)  }%
\end{array}
\right)  =\left(
\begin{array}
[c]{ccc}%
\mathbb{I}_{\mathcal{K}_{OD}} & \mathbb{O} & \mathbb{O}\\
\mathbb{O} & \mathbb{I}_{\mathcal{K}_{D}} & \mathbb{O}\\
\mathbb{O} & \mathbb{O} & \mathbb{I}_{\left[  \ker L_{N}^{1}\right]  ^{c}}%
\end{array}
\right)
\end{align}
which shows that the dual map $\Upsilon^{d}$ to $\Upsilon$, which is defined
by%
\begin{equation}
\left(  \Upsilon^{d}\left(  A\right)  |B\right)  =\left(  A|\Upsilon\left(
B\right)  \right)  \quad\forall B\in\mathcal{B}_{1}\left(  \mathcal{H}%
^{N}\right)  \quad\forall A\in\mathcal{B}\left(  \mathcal{H}^{N}\right)  ,
\end{equation}
is given by
\begin{align}
\Upsilon_{\Gamma}^{d}\left(  \underline{\left(  \left\vert \mathsf{z}%
^{N}\right\rangle _{2}\left\langle \mathsf{z}^{\prime N}\right\vert
\right\vert },\underline{\left(  k^{d}\right\vert },\left(  \underline
{\mathfrak{x}^{d}}\right\vert \right)   &  =\left(  \underline{\left(
\left\vert \mathsf{z}^{N}\right\rangle _{2}\left\langle \mathsf{z}^{\prime
N}\right\vert \right\vert },\underline{\left(  k^{d}\right\vert },\left(
\underline{\mathfrak{x}^{d}}\right\vert \right)  \mathbf{\Upsilon}_{\Gamma
}^{d}\nonumber\\
&  =\left(  \underline{\left(  \left\vert \mathsf{z}^{N}\right\rangle
_{2}\left\langle \mathsf{z}^{\prime N}\right\vert \right\vert },\underline
{\left(  k^{d}\right\vert },\left(  \underline{\mathfrak{x}^{d}}\right\vert
\right)  \mathbf{\Upsilon}_{\Gamma}^{-1}%
\end{align}
where
\begin{equation}
\Upsilon_{\Gamma}\left(  \underline{\left\vert \left\vert \mathsf{z}%
^{N}\right\rangle _{2}\left\langle \mathsf{z}^{\prime N}\right\vert \right)
},\underline{\left\vert k\right)  },\underline{\left\vert \mathfrak{x}\right)
}\right)  =\left(  \underline{\left\vert \left\vert \mathsf{z}^{N}%
\right\rangle _{2}\left\langle \mathsf{z}^{\prime N}\right\vert \right)
},\underline{\left\vert k\right)  },\underline{\left\vert \mathfrak{x}\right)
}\right)  \mathbf{\Upsilon}_{\Gamma}%
\end{equation}
and
\begin{subequations}
\begin{align}
\mathbf{\Upsilon}_{\Gamma}  &  =\left(
\begin{array}
[c]{ccc}%
\mathbb{I}_{\mathcal{K}_{OD}} & \mathbb{O} & \mathbf{\Gamma}_{\mathcal{K}%
_{OD}}\\
\mathbb{O} & \mathbb{I}_{\mathcal{K}_{D}} & \mathbf{\Gamma}_{\mathcal{K}_{D}%
}\\
\mathbb{O} & \mathbb{O} & \mathbb{I}_{\left[  \ker L_{N}^{1}\right]  ^{c}}%
\end{array}
\right) \\
\mathbf{\Upsilon}_{\Gamma}^{-1}  &  \mathbb{=}\left(
\begin{array}
[c]{ccc}%
\mathbb{I}_{\mathcal{K}_{OD}} & \mathbb{O} & -\mathbf{\Gamma}_{\mathcal{K}%
_{OD}}\\
\mathbb{O} & \mathbb{I}_{\mathcal{K}_{D}} & -\mathbf{\Gamma}_{\mathcal{K}_{D}%
}\\
\mathbb{O} & \mathbb{O} & \mathbb{I}_{\left[  \ker L_{N}^{1}\right]  ^{c}}%
\end{array}
\right)
\end{align}


The non canonical representations, $\tilde{\Gamma}\left(  \mathcal{V}\right)
$, of the base of the vector bundle $\mathfrak{B}_{L_{N}^{1}}$ are thus given
by
\end{subequations}
\begin{align}
\mathcal{V}_{\Gamma}  &  =\left\{  \left(  \underline{\left\vert \left\vert
\mathsf{z}^{N}\right\rangle _{2}\left\langle \mathsf{z}^{\prime N}\right\vert
\right)  },\underline{\left\vert k\right)  },\underline{\left\vert
\mathfrak{x}\right)  }\right)  \left(
\begin{array}
[c]{ccc}%
\mathbb{I}_{\mathcal{K}_{OD}} & \mathbb{O} & \mathbf{\Gamma}_{\mathcal{K}%
_{OD}}\\
\mathbb{O} & \mathbb{I}_{\mathcal{K}_{D}} & \mathbf{\Gamma}_{\mathcal{K}_{D}%
}\\
\mathbb{O} & \mathbb{O} & \mathbb{I}_{\left[  \ker L_{N}^{1}\right]  ^{c}}%
\end{array}
\right)  \left(
\begin{array}
[c]{c}%
\mathbf{0}\\
\mathbf{0}\\
\mathfrak{y}%
\end{array}
\right)  ;\mathfrak{y\in}\left[  \ker L_{N}^{1}\right]  ^{c}\right\}
\nonumber\\
&  =\left\{  \left(  \underline{\left\vert \left\vert \mathsf{z}%
^{N}\right\rangle _{2}\left\langle \mathsf{z}^{\prime N}\right\vert \right)
},\underline{\left\vert k\right)  },\underline{\left\vert \mathfrak{x}\right)
}\right)  \left(
\begin{array}
[c]{c}%
\mathbf{\Gamma}_{\mathcal{K}_{OD}}\mathfrak{y}\\
\mathbf{\Gamma}_{\mathcal{K}_{D}}\mathfrak{y}\\
\mathfrak{y}%
\end{array}
\right)  ;\mathfrak{y\in}\left[  \ker L_{N}^{1}\right]  ^{c}\right\}
\nonumber\\
&  =\left\{  \left(  \mathfrak{b}_{0},\cdots,\mathfrak{b}_{N}\right)  \left(
\begin{array}
[c]{c}%
\mathfrak{b}_{0}\\
\vdots\\
\mathfrak{b}_{N}%
\end{array}
\right)  \left(  \underline{\left\vert \left\vert \mathsf{z}^{N}\right\rangle
_{2}\left\langle \mathsf{z}^{\prime N}\right\vert \right)  },\underline
{\left\vert k\right)  },\underline{\left\vert \mathfrak{x}\right)  }\right)
\left(
\begin{array}
[c]{c}%
\mathbf{\Gamma}_{\mathcal{K}_{OD}}\mathfrak{y}\\
\mathbf{\Gamma}_{\mathcal{K}_{D}}\mathfrak{y}\\
\mathfrak{y}%
\end{array}
\right)  ;\mathfrak{y\in}\left[  \ker L_{N}^{1}\right]  ^{c}\right\}
\nonumber\\
&  =\left\{  \left(  \mathfrak{b}_{0},\cdots,\mathfrak{b}_{N}\right)  \left(
\begin{array}
[c]{ccc}%
\mathbb{I} & \mathbb{O} & \mathbb{O}\\
\mathbb{O} & \left(  \mathfrak{b}_{N-1}|\left\vert \underline{k}\right)
\right)  & \left(  \mathfrak{b}_{N}|\left\vert \underline{k}\right)  \right)
\\
\mathbb{O} & \left(  \mathfrak{b}_{N-1}|\underline{\left\vert \mathfrak{x}%
\right)  }\right)  & \left(  \mathfrak{b}_{N}|\underline{\left\vert
\mathfrak{x}\right)  }\right)
\end{array}
\right)  \left(
\begin{array}
[c]{c}%
\mathbf{\Gamma}_{\mathcal{K}_{OD}}\mathfrak{y}\\
\mathbf{\Gamma}_{\mathcal{K}_{D}}\mathfrak{y}\\
\mathfrak{y}%
\end{array}
\right)  ;\mathfrak{y\in}\left[  \ker L_{N}^{1}\right]  ^{c}\right\}
\end{align}


\subsubsection{Non Canonical Basis of Complimentary Subspaces of $\ker
L_{N}^{1}.$}

We can utilize the maps $\tilde{\Gamma}$ introduced in eq. $\left(
\ref{Csubspace}\right)  $ to define a map $\Upsilon:\mathcal{B}_{1}\left(
\mathcal{H}^{N}\right)  \rightarrow\mathcal{B}_{1}\left(  \mathcal{H}%
^{N}\right)  $ to produce bases adapted to the decompositions eq. $\left(
\ref{nonorthog}\right)  $ as
\begin{align}
&  \left\{  \left\{  \left\vert \mathsf{z}^{N}\right\rangle \left\langle
\mathsf{z}^{\prime N}\right\vert \,\forall\mathsf{z}^{N}\overset{2}{\nsim
}\mathsf{z}^{\prime N}\right\}  ,\left\{  k_{\kappa}\,;1\leq\lambda\leq
\infty\right\}  ,\left\{  \tilde{\Gamma}\left(  \mathfrak{x}_{\lambda}\right)
\,;1\leq\lambda\leq\infty\right\}  \right\} \nonumber\\
&  \equiv\left\{  \left\{  \Upsilon\left(  \left\vert \mathsf{z}%
^{N}\right\rangle \left\langle \mathsf{z}^{\prime N}\right\vert \right)
\,\forall\mathsf{z}^{N}\overset{2}{\nsim}\mathsf{z}^{\prime N}\right\}
,\left\{  \Upsilon\left(  k_{\kappa}\right)  \,;1\leq\lambda\leq
\infty\right\}  ,\left\{  \Upsilon\left(  \mathfrak{x}_{\lambda}\right)
\,;1\leq\lambda\leq\infty\right\}  \right\}  \label{noncanbasis}%
\end{align}
where one should note that it is only the bases elements of $\left[  \ker
L_{N}^{1}\right]  ^{c}$ that have been changed. The\emph{\ }dual basis
elements $\left\{  \left(  k_{\kappa}^{d}\right\vert \right\}  $ are now no
longer biorthogonal to $\left\{  \left\vert \tilde{\Gamma}\left(
\mathfrak{x}_{\lambda}\right)  \right)  \right\}  $ as%
\begin{equation}
\tilde{\Gamma}\left(  \mathfrak{x}_{\lambda}\right)  =\Gamma\left(
\mathfrak{x}_{\lambda}\right)  +\mathfrak{x}_{\lambda}.
\end{equation}
and
\begin{equation}
\Gamma\left(  \mathfrak{x}_{\lambda}\right)  =\sum_{\kappa}\left\vert
k_{\kappa}\right)  \left[  \Gamma_{\mathcal{K}_{D}}\right]  _{\kappa\lambda}+%
%TCIMACRO{\dint _{\mathsf{z}^{N}\,\overset{2}{\neq\,}\mathsf{z}^{\prime N}}}%
%BeginExpansion
{\displaystyle\int_{\mathsf{z}^{N}\,\overset{2}{\neq\,}\mathsf{z}^{\prime N}}}
%EndExpansion
\left[  \Gamma_{\mathcal{K}_{OD}}\right]  _{\lambda}\left(  \mathsf{z}%
^{N};\mathsf{z}^{\prime N}\right)  \left\vert \mathsf{z}^{N}\right\rangle
\left\langle \mathsf{z}^{\prime N}\right\vert d\mathsf{z}^{N}d\mathsf{z}%
^{\prime N}%
\end{equation}
However one can construct a basis that is biorthogonal to eq. $\left(
\ref{noncanbasis}\right)  $ by considering
\begin{align}
&  \left(  \Upsilon_{\Gamma}^{d}\left(  \underline{\left(  \left\vert
\mathsf{z}^{N}\right\rangle _{2}\left\langle \mathsf{z}^{\prime N}\right\vert
\right\vert },\underline{\left(  k^{d}\right\vert },\left(  \underline
{\mathfrak{x}^{d}}\right\vert \right)  \right)  \left(  \Upsilon_{\Gamma
}\left(  \underline{\left\vert \left\vert \mathsf{z}^{N}\right\rangle
_{2}\left\langle \mathsf{z}^{\prime N}\right\vert \right)  },\underline
{\left\vert k\right)  },\underline{\left\vert \mathfrak{x}\right)  }\right)
\right)  ^{t}=\nonumber\\
&  \left(  \underline{\left(  \left\vert \mathsf{z}^{N}\right\rangle
_{2}\left\langle \mathsf{z}^{\prime N}\right\vert \right\vert },\underline
{\left(  k^{d}\right\vert },\left(  \underline{\mathfrak{x}^{d}}\right\vert
\right)  \left(
\begin{array}
[c]{ccc}%
\mathbb{I} & \mathbb{O} & -\mathbf{\Gamma}_{\mathcal{K}_{OD}}\\
\mathbb{O} & \mathbb{I} & -\mathbf{\Gamma}_{\mathcal{K}_{D}}\\
\mathbb{O} & \mathbb{O} & \mathbb{I}%
\end{array}
\right)  \left(
\begin{array}
[c]{ccc}%
\mathbb{I} & \mathbb{O} & \mathbf{\Gamma}_{\mathcal{K}_{OD}}\\
\mathbb{O} & \mathbb{I} & \mathbf{\Gamma}_{\mathcal{K}_{D}}\\
\mathbb{O} & \mathbb{O} & \mathbb{I}%
\end{array}
\right)  \left(
\begin{array}
[c]{c}%
\underline{\left\vert \left\vert \mathsf{z}^{N}\right\rangle _{2}\left\langle
\mathsf{z}^{\prime N}\right\vert \right)  }\\
\underline{\left\vert k\right)  }\\
\underline{\left\vert \mathfrak{x}\right)  }%
\end{array}
\right) \nonumber\\
&  \left.  =\right.  \left(  \underline{\left(  \left\vert \mathsf{z}%
^{N}\right\rangle _{2}\left\langle \mathsf{z}^{\prime N}\right\vert
\right\vert },\underline{\left(  k^{d}\right\vert },\left(  \underline
{\mathfrak{x}^{d}}\right\vert \right)  \left(
\begin{array}
[c]{c}%
\underline{\left\vert \left\vert \mathsf{z}^{N}\right\rangle _{2}\left\langle
\mathsf{z}^{\prime N}\right\vert \right)  }\\
\underline{\left\vert k\right)  }\\
\underline{\left\vert \mathfrak{x}\right)  }%
\end{array}
\right)  =\left(
\begin{array}
[c]{ccc}%
\mathbb{I}_{\mathcal{K}_{OD}} & \mathbb{O} & \mathbb{O}\\
\mathbb{O} & \mathbb{I}_{\mathcal{K}_{D}} & \mathbb{O}\\
\mathbb{O} & \mathbb{O} & \mathbb{I}_{\left[  \ker L_{N}^{1}\right]  ^{c}}%
\end{array}
\right)
\end{align}
which shows that the dual map $\Upsilon^{d}$ to $\Upsilon$, which is defined
by%
\begin{equation}
\left(  \Upsilon^{d}\left(  A\right)  |B\right)  =\left(  A|\Upsilon\left(
B\right)  \right)  \quad\forall B\in\mathcal{B}_{1}\left(  \mathcal{H}%
^{N}\right)  \quad\forall A\in\mathcal{B}\left(  \mathcal{H}^{N}\right)  ,
\end{equation}
is given by
\begin{align}
\Upsilon_{\Gamma}^{d}\left(  \underline{\left(  \left\vert \mathsf{z}%
^{N}\right\rangle _{2}\left\langle \mathsf{z}^{\prime N}\right\vert
\right\vert },\underline{\left(  k^{d}\right\vert },\left(  \underline
{\mathfrak{x}^{d}}\right\vert \right)   &  =\left(  \underline{\left(
\left\vert \mathsf{z}^{N}\right\rangle _{2}\left\langle \mathsf{z}^{\prime
N}\right\vert \right\vert },\underline{\left(  k^{d}\right\vert },\left(
\underline{\mathfrak{x}^{d}}\right\vert \right)  \mathbf{\Upsilon}_{\Gamma
}^{d}\nonumber\\
&  =\left(  \underline{\left(  \left\vert \mathsf{z}^{N}\right\rangle
_{2}\left\langle \mathsf{z}^{\prime N}\right\vert \right\vert },\underline
{\left(  k^{d}\right\vert },\left(  \underline{\mathfrak{x}^{d}}\right\vert
\right)  \mathbf{\Upsilon}_{\Gamma}^{-1}%
\end{align}
where
\begin{equation}
\Upsilon_{\Gamma}\left(  \underline{\left\vert \left\vert \mathsf{z}%
^{N}\right\rangle _{2}\left\langle \mathsf{z}^{\prime N}\right\vert \right)
},\underline{\left\vert k\right)  },\underline{\left\vert \mathfrak{x}\right)
}\right)  =\left(  \underline{\left\vert \left\vert \mathsf{z}^{N}%
\right\rangle _{2}\left\langle \mathsf{z}^{\prime N}\right\vert \right)
},\underline{\left\vert k\right)  },\underline{\left\vert \mathfrak{x}\right)
}\right)  \mathbf{\Upsilon}_{\Gamma}%
\end{equation}
and
\begin{subequations}
\begin{align}
\mathbf{\Upsilon}_{\Gamma}  &  =\left(
\begin{array}
[c]{ccc}%
\mathbb{I}_{\mathcal{K}_{OD}} & \mathbb{O} & \mathbf{\Gamma}_{\mathcal{K}%
_{OD}}\\
\mathbb{O} & \mathbb{I}_{\mathcal{K}_{D}} & \mathbf{\Gamma}_{\mathcal{K}_{D}%
}\\
\mathbb{O} & \mathbb{O} & \mathbb{I}_{\left[  \ker L_{N}^{1}\right]  ^{c}}%
\end{array}
\right) \\
\mathbf{\Upsilon}_{\Gamma}^{-1}  &  \mathbb{=}\left(
\begin{array}
[c]{ccc}%
\mathbb{I}_{\mathcal{K}_{OD}} & \mathbb{O} & -\mathbf{\Gamma}_{\mathcal{K}%
_{OD}}\\
\mathbb{O} & \mathbb{I}_{\mathcal{K}_{D}} & -\mathbf{\Gamma}_{\mathcal{K}_{D}%
}\\
\mathbb{O} & \mathbb{O} & \mathbb{I}_{\left[  \ker L_{N}^{1}\right]  ^{c}}%
\end{array}
\right)
\end{align}


The non canonical representations, $\tilde{\Gamma}\left(  \mathcal{V}\right)
$, of the base of the vector bundle $\mathfrak{B}_{L_{N}^{1}}$ are thus given
by
\end{subequations}
\begin{align}
\mathcal{V}_{\Gamma}  &  =\left\{  \left(  \underline{\left\vert \left\vert
\mathsf{z}^{N}\right\rangle _{2}\left\langle \mathsf{z}^{\prime N}\right\vert
\right)  },\underline{\left\vert k\right)  },\underline{\left\vert
\mathfrak{x}\right)  }\right)  \left(
\begin{array}
[c]{ccc}%
\mathbb{I}_{\mathcal{K}_{OD}} & \mathbb{O} & \mathbf{\Gamma}_{\mathcal{K}%
_{OD}}\\
\mathbb{O} & \mathbb{I}_{\mathcal{K}_{D}} & \mathbf{\Gamma}_{\mathcal{K}_{D}%
}\\
\mathbb{O} & \mathbb{O} & \mathbb{I}_{\left[  \ker L_{N}^{1}\right]  ^{c}}%
\end{array}
\right)  \left(
\begin{array}
[c]{c}%
\mathbf{0}\\
\mathbf{0}\\
\mathfrak{y}%
\end{array}
\right)  ;\mathfrak{y\in}\left[  \ker L_{N}^{1}\right]  ^{c}\right\}
\nonumber\\
&  =\left\{  \left(  \underline{\left\vert \left\vert \mathsf{z}%
^{N}\right\rangle _{2}\left\langle \mathsf{z}^{\prime N}\right\vert \right)
},\underline{\left\vert k\right)  },\underline{\left\vert \mathfrak{x}\right)
}\right)  \left(
\begin{array}
[c]{c}%
\mathbf{\Gamma}_{\mathcal{K}_{OD}}\mathfrak{y}\\
\mathbf{\Gamma}_{\mathcal{K}_{D}}\mathfrak{y}\\
\mathfrak{y}%
\end{array}
\right)  ;\mathfrak{y\in}\left[  \ker L_{N}^{1}\right]  ^{c}\right\}
\nonumber\\
&  =\left\{  \left(  \mathfrak{b}_{0},\cdots,\mathfrak{b}_{N}\right)  \left(
\begin{array}
[c]{c}%
\mathfrak{b}_{0}\\
\vdots\\
\mathfrak{b}_{N}%
\end{array}
\right)  \left(  \underline{\left\vert \left\vert \mathsf{z}^{N}\right\rangle
_{2}\left\langle \mathsf{z}^{\prime N}\right\vert \right)  },\underline
{\left\vert k\right)  },\underline{\left\vert \mathfrak{x}\right)  }\right)
\left(
\begin{array}
[c]{c}%
\mathbf{\Gamma}_{\mathcal{K}_{OD}}\mathfrak{y}\\
\mathbf{\Gamma}_{\mathcal{K}_{D}}\mathfrak{y}\\
\mathfrak{y}%
\end{array}
\right)  ;\mathfrak{y\in}\left[  \ker L_{N}^{1}\right]  ^{c}\right\}
\nonumber\\
&  =\left\{  \left(  \mathfrak{b}_{0},\cdots,\mathfrak{b}_{N}\right)  \left(
\begin{array}
[c]{ccc}%
\mathbb{I} & \mathbb{O} & \mathbb{O}\\
\mathbb{O} & \left(  \mathfrak{b}_{N-1}|\left\vert \underline{k}\right)
\right)  & \left(  \mathfrak{b}_{N}|\left\vert \underline{k}\right)  \right)
\\
\mathbb{O} & \left(  \mathfrak{b}_{N-1}|\underline{\left\vert \mathfrak{x}%
\right)  }\right)  & \left(  \mathfrak{b}_{N}|\underline{\left\vert
\mathfrak{x}\right)  }\right)
\end{array}
\right)  \left(
\begin{array}
[c]{c}%
\mathbf{\Gamma}_{\mathcal{K}_{OD}}\mathfrak{y}\\
\mathbf{\Gamma}_{\mathcal{K}_{D}}\mathfrak{y}\\
\mathfrak{y}%
\end{array}
\right)  ;\mathfrak{y\in}\left[  \ker L_{N}^{1}\right]  ^{c}\right\}
\end{align}


\subsubsection{Non Canonical Bases that Contain States \label{noncanstate}}

We concentrate on those representations of the base that are generated by the
set $\mathfrak{P}$ of maps $\left\{  \Gamma\right\}  $ that induce base
representatives and contains a maximal number of states
\begin{equation}
\mathfrak{P}=\left\{  \Gamma|\tilde{\Gamma}\left(  \mathfrak{d}_{1}\right)
\in\mathcal{B}_{1+}\left(  \mathcal{H}^{N}\right)  \,\forall\,\mathfrak{d}%
_{1}\in\ \mathfrak{C}_{1}\right\}
\end{equation}
where $\mathfrak{C}_{1}$ is the set of $\left[  \ker L_{N}^{1}\right]  ^{c}%
$-components of unnormalized states i.e.
\begin{equation}
\mathfrak{C}_{1}=\left\{  \mathfrak{d}_{1}|\mathfrak{d}_{1}\in\left[  \ker
L_{N}^{1}\right]  ^{c}\,\text{such that }\exists\mathfrak{d\in}\left[  \ker
L_{N}^{1}\right]  \text{ so that }\mathfrak{d+d}_{1}\in\mathcal{B}_{1+}\left(
\mathcal{H}^{N}\right)  \right\}
\end{equation}
Thus the cone
\begin{equation}
\mathcal{V}_{\Gamma+}=\mathcal{V}_{\Gamma}\cap\mathcal{B}_{1+}\left(
\mathcal{H}^{N}\right)  \equiv\left\{  \tilde{\Gamma}\left(  \mathfrak{d}%
_{1}\right)  |\mathfrak{d}_{1}\in\ \mathfrak{C}_{1};\tilde{\Gamma}\left(
\mathfrak{d}_{1}\right)  \in\mathcal{B}_{1+}\left(  \mathcal{H}^{N}\right)
\right\}
\end{equation}
is given by
\begin{equation}
\mathcal{V}_{\Gamma+}=\left\{  \left(  \underline{\left\vert \left\vert
\mathsf{z}^{N}\right\rangle _{2}\left\langle \mathsf{z}^{\prime N}\right\vert
\right)  },\underline{\left\vert k\right)  },\underline{\left\vert
\mathfrak{x}\right)  }\right)  \left(
\begin{array}
[c]{ccc}%
\mathbb{I}_{\mathcal{K}_{OD}} & \mathbb{O} & \mathbf{\Gamma}_{\mathcal{K}%
_{OD}}\\
\mathbb{O} & \mathbb{I}_{\mathcal{K}_{D}} & \mathbf{\Gamma}_{\mathcal{K}_{D}%
}\\
\mathbb{O} & \mathbb{O} & \mathbb{I}_{\left[  \ker L_{N}^{1}\right]  ^{c}}%
\end{array}
\right)  \left(
\begin{array}
[c]{c}%
\mathbf{0}\\
\mathbf{0}\\
\mathfrak{d}_{1}%
\end{array}
\right)  ;\mathfrak{d}_{1}\mathfrak{\in C}_{1}\right\}
\end{equation}
We denote a typical element of $\mathcal{V}_{\Gamma+}$ as $D_{0_{F}}^{N}$ thus
for $\mathfrak{d}_{1}\mathfrak{\in C}_{1}$
\begin{align}
D_{0_{F}}^{N}  &  =\Upsilon_{\Gamma}\left(  \mathfrak{d}_{1}\right)
=\tilde{\Gamma}\left(  \mathfrak{d}_{1}\right)  =\left(  \underline{\left\vert
\left\vert \mathsf{z}^{N}\right\rangle _{2}\left\langle \mathsf{z}^{\prime
N}\right\vert \right)  },\underline{\left\vert k\right)  },\underline
{\left\vert \mathfrak{x}\right)  }\right)  \left(
\begin{array}
[c]{c}%
\mathbf{\Gamma}_{\mathcal{K}_{OD}}\mathfrak{d}_{1}\\
\mathbf{\Gamma}_{\mathcal{K}_{D}}\mathfrak{d}_{1}\\
\mathfrak{d}_{1}%
\end{array}
\right) \nonumber\\
&  =\left(  \mathfrak{b}_{0},\cdots,\mathfrak{b}_{N}\right)  \left(
\begin{array}
[c]{ccc}%
\mathbb{I} & \mathbb{O} & \mathbb{O}\\
\mathbb{O} & \left(  \mathfrak{b}_{N-1}|\left\vert \underline{k}\right)
\right)  & \left(  \mathfrak{b}_{N}|\left\vert \underline{k}\right)  \right)
\\
\mathbb{O} & \left(  \mathfrak{b}_{N-1}|\underline{\left\vert \mathfrak{x}%
\right)  }\right)  & \left(  \mathfrak{b}_{N}|\underline{\left\vert
\mathfrak{x}\right)  }\right)
\end{array}
\right)  \left(
\begin{array}
[c]{c}%
\mathbf{\Gamma}_{\mathcal{K}_{OD}}\mathfrak{d}_{1}\\
\mathbf{\Gamma}_{\mathcal{K}_{D}}\mathfrak{d}_{1}\\
\mathfrak{d}_{1}%
\end{array}
\right) \nonumber\\
&  =\left(  \mathfrak{b}_{0},\cdots,\mathfrak{b}_{N}\right)  \left(
\begin{array}
[c]{c}%
D_{0_{F}\mathfrak{B}_{0}}^{N}\\
\vdots\\
D_{0_{F}\mathfrak{B}_{N}}^{N}%
\end{array}
\right)
\end{align}
recall that the domain of $\Upsilon_{\Gamma}$ is all of $\mathcal{B}%
_{1}\left(  \mathcal{H}^{N}\right)  $, while that of $\tilde{\Gamma}$ is
$\left[  \ker L_{N}^{1}\right]  ^{c}$ and
\begin{equation}
D_{0_{F}\mathfrak{B}_{m}}^{N}=%
%TCIMACRO{\dint \limits_{\mathsf{z}^{\prime N}\,\overset{m}{=}\,\mathsf{z}^{N}%
%}}%
%BeginExpansion
{\displaystyle\int\limits_{\mathsf{z}^{\prime N}\,\overset{m}{=}%
\,\mathsf{z}^{N}}}
%EndExpansion
D_{0_{F}}^{N}\left(  \mathsf{z}^{\prime N},\mathsf{z}^{N}\right)  \left\vert
\mathsf{z}^{\prime N}\right\rangle \left\langle \mathsf{z}^{N}\right\vert
d\mathsf{z}^{N}d\mathsf{z}^{\prime N}%
\end{equation}
so in particular%
\begin{align}
&  D_{0_{F}\mathfrak{B}_{N}}^{N}=%
%TCIMACRO{\dint }%
%BeginExpansion
{\displaystyle\int}
%EndExpansion
D_{0_{F}}^{N}\left(  \mathsf{z}^{N},\mathsf{z}^{N}\right)  \left\vert
\mathsf{z}^{N}\right\rangle \left\langle \mathsf{z}^{N}\right\vert
d\mathsf{z}^{N}\nonumber\\
&  \qquad\quad\left.  =\right.  \left(  \mathfrak{b}_{N}|\underline{\left\vert
\mathfrak{x}\right)  }\right)  \mathbf{\Gamma}_{\mathcal{K}_{D}}%
\mathfrak{d}_{1}+\left(  \mathfrak{b}_{N}|\underline{\left\vert \mathfrak{x}%
\right)  }\right)  \mathfrak{d}_{1}\nonumber\\
&  \qquad\quad\left.  =\right.  \left(  \mathfrak{b}_{N}|\underline{\left\vert
k\right)  }\right)  \left(  \underline{\left(  k^{d}\right\vert }%
|\Gamma\left(  \mathfrak{d}_{1}\right)  \right)  +\left(  \mathfrak{b}%
_{N}|\mathfrak{d}_{1}\right) \nonumber\\
&  \left.  =\right.
%TCIMACRO{\dsum \limits_{1\leq\nu\leq\infty}}%
%BeginExpansion
{\displaystyle\sum\limits_{1\leq\nu\leq\infty}}
%EndExpansion%
%TCIMACRO{\dint }%
%BeginExpansion
{\displaystyle\int}
%EndExpansion
k_{\nu}\left(  \mathsf{z}^{N},\mathsf{z}^{N}\right)  \Gamma\left[
\mathfrak{d}_{1}\right]  \left(  \mathsf{z}^{N},\mathsf{z}^{N}\right)
\left\vert \mathsf{z}^{N}\right\rangle \left\langle \mathsf{z}^{N}\right\vert
d\mathsf{z}^{N}+\int\mathfrak{d}_{1}\left(  \mathsf{z}^{N},\mathsf{z}%
^{N}\right)  \left\vert \mathsf{z}^{N}\right\rangle \left\langle
\mathsf{z}^{N}\right\vert d\mathsf{z}^{N}%
\end{align}
so that
\begin{equation}
D_{0_{F}}^{N}\left(  \mathsf{z}^{N},\mathsf{z}^{N}\right)  =D_{0_{F}%
\mathfrak{B}_{N}}^{N}\left(  \mathsf{z}^{N},\mathsf{z}^{N}\right)  =%
%TCIMACRO{\dsum \limits_{1\leq\nu\leq\infty}}%
%BeginExpansion
{\displaystyle\sum\limits_{1\leq\nu\leq\infty}}
%EndExpansion
k_{\nu}\left(  \mathsf{z}^{N},\mathsf{z}^{N}\right)  \Gamma\left[
\mathfrak{d}_{1}\right]  \left(  \mathsf{z}^{N},\mathsf{z}^{N}\right)
+\mathfrak{d}_{1}\left(  \mathsf{z}^{N},\mathsf{z}^{N}\right)
\end{equation}
while
\begin{align}
D_{0_{F}\mathfrak{B}_{N-1}}^{N}  &  =%
%TCIMACRO{\dint \limits_{\mathsf{z}^{\prime N}\,\overset{N-1}{=}\,\mathsf{z}%
%^{N}}}%
%BeginExpansion
{\displaystyle\int\limits_{\mathsf{z}^{\prime N}\,\overset{N-1}{=}%
\,\mathsf{z}^{N}}}
%EndExpansion
D_{0_{F}}^{N}\left(  \mathsf{z}^{\prime N},\mathsf{z}^{N}\right)  \left\vert
\mathsf{z}^{\prime N}\right\rangle \left\langle \mathsf{z}^{N}\right\vert
d\mathsf{z}^{N}d\mathsf{z}^{\prime N}\nonumber\\
&  =\left(  \mathfrak{b}_{N-1}|\underline{\left\vert \mathfrak{x}\right)
}\right)  \mathbf{\Gamma}_{\mathcal{K}_{D}}\mathfrak{d}_{1}+\left(
\mathfrak{b}_{N-1}|\underline{\left\vert \mathfrak{x}\right)  }\right)
\mathfrak{d}_{1}\nonumber\\
&  =\left(  \mathfrak{b}_{N-1}|\underline{\left\vert k\right)  }\right)
\left(  \underline{\left(  k^{d}\right\vert }|\Gamma\left(  \mathfrak{d}%
_{1}\right)  \right)  +\left(  \mathfrak{b}_{N-1}|\mathfrak{d}_{1}\right)
\nonumber\\
&  =%
%TCIMACRO{\dsum \limits_{1\leq\nu\leq\infty}}%
%BeginExpansion
{\displaystyle\sum\limits_{1\leq\nu\leq\infty}}
%EndExpansion%
%TCIMACRO{\dint \limits_{\mathsf{z}^{\prime N}\,\overset{N-1}{=}\,\mathsf{z}%
%^{N}}}%
%BeginExpansion
{\displaystyle\int\limits_{\mathsf{z}^{\prime N}\,\overset{N-1}{=}%
\,\mathsf{z}^{N}}}
%EndExpansion
k_{\nu}\left(  \mathsf{z}^{N},\mathsf{z}^{\prime N}\right)  \Gamma\left[
\mathfrak{d}_{1}\right]  \left(  \mathsf{z}^{N},\mathsf{z}^{\prime N}\right)
\left\vert \mathsf{z}^{N}\right\rangle \left\langle \mathsf{z}^{\prime
N}\right\vert d\mathsf{z}^{N}\nonumber\\
&  \qquad\qquad\qquad\qquad\qquad\qquad+%
%TCIMACRO{\dint \limits_{\mathsf{z}^{\prime N}\,\overset{N-1}{=}\,\mathsf{z}%
%^{N}}}%
%BeginExpansion
{\displaystyle\int\limits_{\mathsf{z}^{\prime N}\,\overset{N-1}{=}%
\,\mathsf{z}^{N}}}
%EndExpansion
\mathfrak{d}_{1}\left(  \mathsf{z}^{N},\mathsf{z}^{\prime N}\right)
\left\vert \mathsf{z}^{N}\right\rangle \left\langle \mathsf{z}^{\prime
N}\right\vert d\mathsf{z}^{N}d\mathsf{z}^{\prime N}%
\end{align}
so that
\begin{equation}
D_{0_{F}}^{N}\left(  \mathsf{z}^{\prime N},\mathsf{z}^{N}\right)  =%
%TCIMACRO{\dsum \limits_{1\leq\nu\leq\infty}}%
%BeginExpansion
{\displaystyle\sum\limits_{1\leq\nu\leq\infty}}
%EndExpansion
k_{\nu}\left(  \mathsf{z}^{N},\mathsf{z}^{\prime N}\right)  \Gamma\left[
\mathfrak{d}_{1}\right]  \left(  \mathsf{z}^{N},\mathsf{z}^{\prime N}\right)
+\mathfrak{d}_{1}\left(  \mathsf{z}^{N},\mathsf{z}^{\prime N}\right)
;\quad\mathsf{z}^{\prime N}\overset{N-1}{=}\mathsf{z}^{N}%
\end{equation}
where we have used the notation $\Gamma\left[  \mathfrak{d}_{1}\right]
\left(  \mathsf{z}^{N},\mathsf{z}^{\prime N}\right)  $ to distinguish between
arguments that are representation dependent and those that are not.


\end{document}